% !TeX program = lualatex
\documentclass[12pt,aspectratio=169]{beamer}
\usepackage[utf8]{inputenc}
\usepackage[T1]{fontenc}
\usepackage{amsmath,amsfonts,amssymb}
\usepackage{graphicx}
\usepackage{xcolor}
\usepackage{hyperref}
\usepackage{anyfontsize}
\usepackage{lmodern}
\usepackage{longtable}
\usepackage{array}
\usepackage{booktabs}

% ====== EXTREME FONT REDUCTION ======
\renewcommand{\tiny}{\fontsize{4}{5}\selectfont}
\renewcommand{\scriptsize}{\fontsize{5}{6}\selectfont}
\renewcommand{\footnotesize}{\fontsize{6}{7}\selectfont}
\renewcommand{\small}{\fontsize{7}{8}\selectfont}
\renewcommand{\normalsize}{\fontsize{8}{9}\selectfont}
\renewcommand{\large}{\fontsize{9}{10}\selectfont}
\renewcommand{\Large}{\fontsize{10}{11}\selectfont}
\renewcommand{\LARGE}{\fontsize{11}{13}\selectfont}
\renewcommand{\huge}{\fontsize{12}{14}\selectfont}
\renewcommand{\Huge}{\fontsize{14}{17}\selectfont}

% ====== COLOR THEME ======
\definecolor{redcorrect}{RGB}{200,30,30}
\definecolor{bluecorrect}{RGB}{30,80,200}
\definecolor{greencheck}{RGB}{0,150,50}
\definecolor{lightgray}{RGB}{245,245,245}
\definecolor{darkgray}{RGB}{50,50,50}
\definecolor{softyellow}{RGB}{255,248,200}
\definecolor{steppurple}{RGB}{150,50,200}
\definecolor{deepblue}{RGB}{20,60,150}
\definecolor{darkred}{RGB}{150,20,20}
\definecolor{orangewarning}{RGB}{255,140,0}

\setbeamercolor{title}{fg=white,bg=redcorrect}
\setbeamercolor{frametitle}{fg=white,bg=redcorrect}
\setbeamercolor{block title}{fg=white,bg=bluecorrect}
\setbeamercolor{block body}{fg=darkgray,bg=lightgray}
\setbeamercolor{normal text}{fg=darkgray}
\usecolortheme{default}

% ====== CUSTOM COMMANDS ======
\newcommand{\correct}[1]{\textcolor{greencheck}{\textbf{\checkmark}} #1}
\newcommand{\keypoint}[1]{\textcolor{deepblue}{\textbf{#1}}}
\newcommand{\hardconcept}[1]{\textcolor{darkred}{\textbf{#1}}}
\newcommand{\tipbox}[1]{\fcolorbox{bluecorrect}{softyellow}{\parbox{0.88\textwidth}{\centering #1}}}
\newcommand{\stepbox}[1]{%
	\begin{center}
		\fcolorbox{darkgray}{white}{%
			\parbox{0.90\textwidth}{\vspace{0.05cm} #1 \vspace{0.05cm}}%
		}
	\end{center}
}
\newcommand{\wrongbox}[1]{\fcolorbox{redcorrect}{red!10}{\parbox{0.88\textwidth}{\centering \textcolor{redcorrect}{\textbf{\times}} #1}}}
\newcommand{\highlight}[1]{\textcolor{redcorrect}{\textbf{#1}}}
\newcommand{\bluetext}[1]{\textcolor{bluecorrect}{\textbf{#1}}}
\newcommand{\stagebox}[1]{\fcolorbox{steppurple}{white!5}{\parbox{0.88\textwidth}{\centering #1}}}

\begin{document}
	
	% ================================================================
	% SLIDE 1: TITLE
	% ================================================================
	\begin{frame}
		\centering
		\vspace{0.8cm}
		{\fontsize{16}{19}\selectfont \textbf{Domain A: Understanding Financial Crime}}\\[0.15cm]
		{\fontsize{12}{14}\selectfont Harder/Expanded Q\&A Review}\\[0.3cm]
		{\fontsize{9}{11}\selectfont \textit{Detailed Regulatory Principles \& Complex Scenarios}}\\[0.3cm]
		{\fontsize{8}{10}\selectfont Compliance Training Department — Advanced Module}
		\vspace{0.2cm}
		\rule{4cm}{0.5pt}
	\end{frame}
	
	% ================================================================
	% SLIDE 2: AGENDA
	% ================================================================
	\begin{frame}
		\frametitle{\fontsize{11}{13}\selectfont Agenda — Domain A Advanced Topics}
		\scriptsize
		\begin{longtable}{p{0.45\textwidth} p{0.50\textwidth}}
			\toprule
			\keypoint{Topic} & \keypoint{Questions} \\
			\midrule
			Definitions of AML, CFT, Sanctions, Fraud, ABC, Tax Evasion & Q1-Q6 \\
			Financial Systems and Consequences of ML & Q7-Q10 \\
			Predicate Crimes & Q11-Q14 \\
			Manifestation in Banking Sector & Q15-Q20 \\
			Banking Product Risks & Q21-Q24 \\
			High-Risk Sectors & Q25-Q28 \\
			PEP and High-Risk Customer Risks & Q29-Q32 \\
			Other Financial Sector Risks & Q33-Q36 \\
			MSB, PSP, E-Commerce Risks & Q37-Q40 \\
			Insurance Product Risks & Q41-Q44 \\
			VASP and Cryptoasset Risks & Q45-Q48 \\
			Gaming and Gambling Risks & Q49-Q52 \\
			Real Estate Risks & Q53-Q56 \\
			Gatekeeper Risks & Q57-Q60 \\
			TCSP Risks & Q61-Q64 \\
			Accountant Risks & Q65-Q68 \\
			Other High-Risk Sectors & Q69-Q72 \\
			\bottomrule
		\end{longtable}
	\end{frame}
	
	% ================================================================
	% SLIDE 3: Q1 — Money Laundering Definition (EXPANDED)
	% ================================================================
	\begin{frame}
		\frametitle{\fontsize{11}{13}\selectfont Q1: Money Laundering — Complex Scenario}
		\begin{block}{\fontsize{8}{10}\selectfont Topic: Definitions — Advanced Application}
			\scriptsize
			A compliance officer is training new staff on foundational concepts. One trainee asks: \textit{"Isn't money laundering just hiding money from authorities?"} Which response best clarifies the true definition and scope, and addresses the specific limitations of the trainee's understanding?
		\end{block}
		\vspace{0.05cm}
		\stepbox{\scriptsize
			\keypoint{Correct Answer:} Money laundering is the process of making illicit funds appear legitimate by concealing their true origin through a series of transactions, typically involving placement, layering, and integration. "Hiding money" is an oversimplification because the ultimate goal is to create a \textbf{legitimate appearance} so funds can be used without raising suspicion. Hiding alone does not enable use—integration does.
		}
		\vspace{0.05cm}
		\scriptsize
		\hardconcept{Core Concept — Expanded:} The three stages are:
		\par \textbf{Placement}: Introducing illicit funds into the financial system (e.g., structuring deposits).
		\par \textbf{Layering}: Moving funds through complex transactions to obscure origin (e.g., wire transfers, shell companies).
		\par \textbf{Integration}: Returning "cleaned" funds to the legitimate economy (e.g., real estate purchases, investments).
		\vspace{0.03cm}
		\wrongbox{\scriptsize \textbf{Common Trap:} "Money laundering is simply hiding money" — INCORRECT. Hiding is only placement. The \textbf{integration} stage is what makes laundering effective. Without integration, funds remain unusable.}
		\vspace{0.03cm}
		\tipbox{\scriptsize \textbf{Key Insight:} \textit{FATF defines ML as "the processing of criminal proceeds to disguise their illegal origin."} The key is \textbf{legitimization}, not just concealment.}
	\end{frame}
	
	% ================================================================
	% SLIDE 4: Q2 — Terrorist Financing (EXPANDED)
	% ================================================================
	\begin{frame}
		\frametitle{\fontsize{11}{13}\selectfont Q2: Terrorist Financing — Complex Scenario}
		\begin{block}{\fontsize{8}{10}\selectfont Topic: Definitions — CFT Advanced}
			\scriptsize
			An investigator is reviewing a case where a charity organization collected legitimate donations, but some funds were diverted to support extremist groups. The investigator states: \textit{"This isn't money laundering because the funds came from legitimate sources."} How should this be addressed, and what key distinction is the investigator missing?
		\end{block}
		\vspace{0.05cm}
		\stepbox{\scriptsize
			\keypoint{Correct Answer:} Terrorist financing is defined as the provision or collection of funds, by any means, directly or indirectly, with the intention that they should be used to carry out terrorist acts. Unlike money laundering, where funds are \textbf{always illicit}, terrorist financing can involve \textbf{legitimate funds diverted for illegal purposes}. The investigator's error is confusing the \textit{source} of funds with the \textit{intent/purpose} of the financing.
		}
		\vspace{0.05cm}
		\scriptsize
		\hardconcept{Core Concept — Expanded:} Key differences:
		\par \textbf{Source}: ML funds are illicit; TF funds may be legitimate (charity, salary, small donations).
		\par \textbf{Amount}: ML often involves large sums; TF can involve small amounts (under reporting thresholds).
		\par \textbf{Intent}: ML is about disguising origin; TF is about \textbf{purpose} — intent to support terrorism.
		\par \textbf{Destination}: ML focuses on where money came from; TF focuses on \textbf{where money is going}.
		\vspace{0.03cm}
		\wrongbox{\scriptsize \textbf{Common Trap:} "Terrorist financing always involves illegal money" — INCORRECT. TF funds can come from \textbf{legitimate sources} (e.g., employment, charity) but are used for illegal purposes. The \textbf{intent} is the defining element.}
		\vspace{0.03cm}
		\tipbox{\scriptsize \textbf{Key Insight:} Under \textit{FATF Recommendation 5}, countries must criminalize terrorist financing regardless of whether the funds are legitimate or illicit. \textbf{Purpose, not source, is determinative.}}
	\end{frame}
	
	% ================================================================
	% SLIDE 5: Q3 — Sanctions (EXPANDED)
	% ================================================================
	\begin{frame}
		\frametitle{\fontsize{11}{13}\selectfont Q3: Sanctions — Complex Scenario}
		\begin{block}{\fontsize{8}{10}\selectfont Topic: Definitions — Sanctions Advanced}
			\scriptsize
			A compliance officer at an international bank is reviewing a transaction involving a company incorporated in Country A but with ultimate beneficial owners who are nationals of Country B, which is subject to comprehensive sanctions. The transaction involves a payment for goods that are not on any restricted list. The officer states: \textit{"Since the goods aren't restricted, we can process this."} Why is this reasoning flawed?
		\end{block}
		\vspace{0.05cm}
		\stepbox{\scriptsize
			\keypoint{Correct Answer:} Sanctions are legal restrictions imposed by governments or international bodies against designated countries, individuals, or entities to achieve foreign policy or national security objectives. Sanctions can be \textbf{comprehensive} (blocking all transactions with a country) or \textbf{targeted} (blocking specific individuals/entities). The officer's reasoning fails to consider \textbf{ownership/control} — if the beneficial owners are sanctioned nationals, the transaction may be prohibited regardless of the goods' classification.
		}
		\vspace{0.05cm}
		\scriptsize
		\hardconcept{Core Concept — Expanded:} Types of sanctions:
		\par \textbf{Country/Regime Sanctions}: Comprehensive restrictions (e.g., Iran, North Korea).
		\par \textbf{Targeted Sanctions}: Against specific individuals or entities (e.g., OFAC SDN List).
		\par \textbf{Sectoral Sanctions}: Restrictions on specific industries (e.g., Russian energy sector).
		\par \textbf{Asset Freezes}: Freezing assets of designated persons.
		\par \textbf{Ownership/Control Rules}: Many sanctions apply to entities \textit{owned or controlled} by designated persons (e.g., OFAC's 50\% rule).
		\vspace{0.03cm}
		\wrongbox{\scriptsize \textbf{Common Trap:} "Sanctions are optional recommendations" — INCORRECT. Sanctions are \textbf{legally binding} and violations carry severe penalties (civil fines, criminal prosecution).}
		\vspace{0.03cm}
		\tipbox{\scriptsize \textbf{Key Insight:} Sanctions are \textbf{strict liability} — no intent is required for a violation. \textit{OFAC enforces strict liability} — even inadvertent violations can result in fines.}
	\end{frame}
	
	% ================================================================
	% SLIDE 6: Q4 — Fraud (EXPANDED)
	% ================================================================
	\begin{frame}
		\frametitle{\fontsize{11}{13}\selectfont Q4: Fraud — Complex Scenario}
		\begin{block}{\fontsize{8}{10}\selectfont Topic: Definitions — Fraud Advanced}
			\scriptsize
			A bank's fraud team investigates a pattern of new account openings where applicants use slightly altered personal information (address variations, different phone numbers) but the same Social Security number variations. One analyst suggests: \textit{"This is just identity theft — the victims are real people."} Another disagrees: \textit{"This is synthetic identity fraud."} Which analyst is correct, and why?
		\end{block}
		\vspace{0.05cm}
		\stepbox{\scriptsize
			\keypoint{Correct Answer:} Financial fraud is intentional deception or misrepresentation made for personal gain or to damage another individual, involving financial transactions or assets. The second analyst is correct — this is \textbf{synthetic identity fraud}, not traditional identity theft. Synthetic identity fraud involves creating a \textbf{new identity} using a combination of real and fabricated information (e.g., real SSN with fake name/address). Traditional identity theft uses \textit{complete} stolen identity information. The use of \textit{varied} address and phone information with the same SSN indicates identity \textit{manufacturing}, not theft.
		}
		\vspace{0.05cm}
		\scriptsize
		\hardconcept{Core Concept — Expanded:} Common fraud types:
		\par \textbf{Identity Theft}: Using stolen personal information of a real person.
		\par \textbf{Synthetic Identity Fraud}: Creating fake identities using real + fabricated data.
		\par \textbf{Investment Fraud}: False investment opportunities (Ponzi schemes).
		\par \textbf{Payment Fraud}: Unauthorized transactions, card fraud.
		\par \textbf{Insurance Fraud}: False claims, inflated damages.
		\par \textbf{Application Fraud}: Using false information in loan/credit applications.
		\vspace{0.03cm}
		\wrongbox{\scriptsize \textbf{Common Trap:} "All identity fraud is identity theft" — INCORRECT. Synthetic identity fraud is \textbf{different} because it creates a \textbf{new identity} rather than stealing an existing one.}
		\vspace{0.03cm}
		\tipbox{\scriptsize \textbf{Key Insight:} Fraud and money laundering often overlap — fraud generates proceeds that need laundering. \textit{The proceeds of fraud are a predicate offense for money laundering under FATF standards.}}
	\end{frame}
	
	% ================================================================
	% SLIDE 7: Q5 — Anti-Bribery and Corruption (EXPANDED)
	% ================================================================
	\begin{frame}
		\frametitle{\fontsize{11}{13}\selectfont Q5: Anti-Bribery — Complex Scenario}
		\begin{block}{\fontsize{8}{10}\selectfont Topic: Definitions — ABC Advanced}
			\scriptsize
			A multinational corporation is expanding operations into multiple jurisdictions. The compliance team is reviewing facilitation payment policies. One jurisdiction allows facilitation payments under local law but the company is incorporated in the UK and also operates in the US. The regional manager argues: \textit{"Since we are in a jurisdiction that allows facilitation payments, we can make them."} The compliance officer disagrees. Who is correct, and why?
		\end{block}
		\vspace{0.05cm}
		\stepbox{\scriptsize
			\keypoint{Correct Answer:} Bribery is the offering, giving, receiving, or soliciting of anything of value to influence an action or decision. The compliance officer is correct. Under the \textbf{UK Bribery Act 2010}, facilitation payments are \textbf{prohibited} regardless of local law. The UK Bribery Act has \textbf{extraterritorial reach} and applies to UK-based companies operating anywhere in the world. While the US Foreign Corrupt Practices Act (FCPA) allows facilitation payments (small payments to speed routine government actions), the UK Act does not. The company must comply with the \textit{stricter} UK requirements.
		}
		\vspace{0.05cm}
		\scriptsize
		\hardconcept{Core Concept — Expanded:} Key distinctions:
		\par \textbf{Bribery}: Paying for influence (illegal under both FCPA and UK Act).
		\par \textbf{Facilitation Payment}: Small payment to speed routine action (allowed under FCPA; \textbf{prohibited} under UK Bribery Act).
		\par \textbf{Hospitality}: Legitimate business entertainment (must be proportionate, transparent).
		\par \textbf{Corruption}: Abuse of entrusted power for private gain.
		\vspace{0.03cm}
		\wrongbox{\scriptsize \textbf{Common Trap:} "All payments to officials are bribery" — INCORRECT. Facilitation payments are treated differently under different laws. \textit{The UK Bribery Act is stricter than the FCPA.}}
		\vspace{0.03cm}
		\tipbox{\scriptsize \textbf{Key Insight:} \textit{The UK Bribery Act is considered the gold standard} for anti-bribery legislation due to its extraterritorial reach, strict prohibition of facilitation payments, and the corporate offense of failing to prevent bribery.}
	\end{frame}
	
	% ================================================================
	% SLIDE 8: Q6 — Tax Evasion (EXPANDED)
	% ================================================================
	\begin{frame}
		\frametitle{\fontsize{11}{13}\selectfont Q6: Tax Evasion — Complex Scenario}
		\begin{block}{\fontsize{8}{10}\selectfont Topic: Definitions — Tax Evasion Advanced}
			\scriptsize
			A wealthy client moves funds to an offshore jurisdiction with no income tax, structures ownership through a trust, and declares only the minimum required income to tax authorities. The client argues: \textit{"This is legitimate tax avoidance — I'm using legal structures to minimize my tax burden."} The bank's compliance officer suspects tax evasion. What distinction determines whether this is avoidance (legal) or evasion (illegal)?
		\end{block}
		\vspace{0.05cm}
		\stepbox{\scriptsize
			\keypoint{Correct Answer:} Tax evasion is the illegal act of deliberately not paying taxes that are legally owed, typically through concealment or misrepresentation. Tax avoidance is legal — using legitimate methods to reduce tax liability. The key distinction is \textbf{disclosure and truthfulness}. If the client is \textit{concealing} the existence of the offshore assets or \textit{misrepresenting} their value/ownership, it is evasion. If the client is \textit{fully disclosing} all structures and transactions to tax authorities, it may be avoidance. The use of trusts and offshore structures alone is not determinative — it's whether the client is \textbf{honest} about their existence and operations.
		}
		\vspace{0.05cm}
		\scriptsize
		\hardconcept{Core Concept — Expanded:} Distinction:
		\par \textbf{Tax Evasion}: Illegal — concealment, misrepresentation, false statements, willful failure to report.
		\par \textbf{Tax Avoidance}: Legal — using legitimate methods to reduce tax liability (e.g., tax treaties, deductions, exemptions).
		\par \textbf{Aggressive Tax Avoidance}: Legal but potentially questionable schemes (may trigger anti-abuse rules).
		\vspace{0.03cm}
		\wrongbox{\scriptsize \textbf{Common Trap:} "Tax evasion and tax avoidance are the same" — INCORRECT. Evasion is \textbf{illegal}; avoidance is \textbf{legal}. The difference is \textit{honesty and disclosure}.}
		\vspace{0.03cm}
		\tipbox{\scriptsize \textbf{Key Insight:} Tax evasion is a \textbf{predicate offense} for money laundering in many jurisdictions (FATF Recommendation 3). \textit{The proceeds of tax evasion must be laundered to be usable.}}
	\end{frame}
	
	% ================================================================
	% SLIDE 9: Q7 — Economic Consequences (EXPANDED)
	% ================================================================
	\begin{frame}
		\frametitle{\fontsize{11}{13}\selectfont Q7: Economic Consequences of ML}
		\begin{block}{\fontsize{8}{10}\selectfont Topic: Financial Systems — Economic Impact}
			\scriptsize
			A developing country has experienced a significant increase in suspicious transaction reporting. The central bank is concerned about the impact on economic stability. What is the most significant economic consequence of unchecked money laundering?
		\end{block}
		\vspace{0.05cm}
		\stepbox{\scriptsize
			\keypoint{Correct Answer:} Distortion of economic statistics and misallocation of resources, undermining legitimate economic activity. This leads to:
			\par \textbf{GDP Distortion}: Economic data becomes unreliable, affecting policy decisions.
			\par \textbf{Misallocation of Capital}: Capital flows to criminal enterprises rather than productive investment.
			\par \textbf{Increased Volatility}: Sudden capital movements destabilize financial markets.
			\par \textbf{Reduced Foreign Investment}: Countries perceived as high-risk lose foreign direct investment.
			\par \textbf{Asset Price Distortion}: Laundered money inflates real estate, art, and luxury goods markets.
		}
		\vspace{0.03cm}
		\tipbox{\scriptsize \textbf{Key Insight:} \textit{The IMF estimates that money laundering represents 2-5\% of global GDP.} This is a massive distortion of legitimate economic activity.}
	\end{frame}
	
	% ================================================================
	% SLIDE 10: Q8 — Reputational Consequences (EXPANDED)
	% ================================================================
	\begin{frame}
		\frametitle{\fontsize{11}{13}\selectfont Q8: Reputational Consequences}
		\begin{block}{\fontsize{8}{10}\selectfont Topic: Financial Systems — Reputational Impact}
			\scriptsize
			A global bank is fined \$2 billion for AML control failures and publicly named in a regulatory enforcement action. The bank's CEO states: \textit{"The fine is manageable — we can absorb it."} The compliance officer disagrees. Why might the CEO's assessment be short-sighted?
		\end{block}
		\vspace{0.05cm}
		\stepbox{\scriptsize
			\keypoint{Correct Answer:} The reputational consequence is often \textbf{more severe} than the financial penalty. Reputational damage includes:
			\par \textbf{Trust Erosion}: Customers lose confidence and withdraw business.
			\par \textbf{Business Loss}: Correspondent banks may terminate relationships.
			\par \textbf{Brand Damage}: Negative media coverage can last for years.
			\par \textbf{Regulatory Scrutiny}: Increased supervision, more frequent exams.
			\par \textbf{Investor Flight}: Stock price decline, difficulty raising capital.
			\par \textbf{Business Opportunity Loss}: Deals lost to competitors.
			The \$2 billion fine may be \textit{less} damaging than the long-term loss of customer and investor trust.
		}
		\vspace{0.03cm}
		\tipbox{\scriptsize \textbf{Key Insight:} \textit{Reputational damage can be permanent.} Banks have been forced out of markets entirely after enforcement actions.}
	\end{frame}
	
	% ================================================================
	% SLIDE 11: Q9 — Social Consequences (EXPANDED)
	% ================================================================
	\begin{frame}
		\frametitle{\fontsize{11}{13}\selectfont Q9: Social Consequences of ML}
		\begin{block}{\fontsize{8}{10}\selectfont Topic: Financial Systems — Social Impact}
			\scriptsize
			An NGO is studying the impact of money laundering on developing countries. They find that countries with significant money laundering activity also have higher levels of organized crime and corruption. What is the most significant social consequence?
		\end{block}
		\vspace{0.05cm}
		\stepbox{\scriptsize
			\keypoint{Correct Answer:} Increased crime, corruption, and social inequality as illicit funds fuel criminal enterprises. Specific social impacts:
			\par \textbf{Crime Increase}: Laundered funds fuel drug trafficking, human trafficking, and organized crime.
			\par \textbf{Corruption}: Criminal influence over political and judicial institutions.
			\par \textbf{Inequality}: Illicit wealth concentration widens income gaps.
			\par \textbf{Weak Rule of Law}: Undermines legal and democratic institutions.
			\par \textbf{Erosion of Trust}: Citizens lose faith in government and financial systems.
			\par \textbf{Human Rights Violations}: Fueled by trafficking and corruption proceeds.
		}
		\vspace{0.03cm}
		\tipbox{\scriptsize \textbf{Key Insight:} Money laundering is not a victimless crime — \textit{it fuels organized crime, corruption, and human rights abuses}.}
	\end{frame}
	
	% ================================================================
	% SLIDE 12: Q10 — Regulatory Penalties (EXPANDED)
	% ================================================================
	\begin{frame}
		\frametitle{\fontsize{11}{13}\selectfont Q10: Regulatory Penalties}
		\begin{block}{\fontsize{8}{10}\selectfont Topic: Financial Systems — Enforcement}
			\scriptsize
			A compliance officer is advising the board on potential consequences of AML control failures. The board asks: \textit{"Worst case, what could happen?"} What should the compliance officer explain?
		\end{block}
		\vspace{0.05cm}
		\stepbox{\scriptsize
			\keypoint{Correct Answer:} Significant financial fines, criminal prosecution of individuals, and revocation of licenses. Under FATF Recommendation 35, sanctions must be effective, proportionate, and dissuasive:
			\par \textbf{Civil Penalties}: Financial fines (e.g., US BSA penalties up to \$1M+ per violation).
			\par \textbf{Criminal Penalties}: Imprisonment for individuals (e.g., up to 20 years in US).
			\par \textbf{Administrative Penalties}: License suspension or revocation.
			\par \textbf{Personal Liability}: Compliance officers and senior management can be held personally liable.
			\par \textbf{De-risking}: Correspondent banking relationships may be terminated.
			\par \textbf{Supervisory Action}: Increased supervision, remediation requirements.
			The \textit{worst case} can include criminal prosecution of executives and closure of the institution.
		}
		\vspace{0.03cm}
		\tipbox{\scriptsize \textbf{Key Insight:} Under \textit{FATF Recommendation 35}, sanctions must be effective, proportionate, and dissuasive. \textbf{Personal liability} is increasingly common.}
	\end{frame}
	
	% ================================================================
	% SLIDE 13: Q11 — Predicate Offenses (EXPANDED)
	% ================================================================
	\begin{frame}
		\frametitle{\fontsize{11}{13}\selectfont Q11: Predicate Offenses}
		\begin{block}{\fontsize{8}{10}\selectfont Topic: Predicate Crimes — Advanced}
			\scriptsize
			A jurisdiction is updating its AML legislation to align with FATF standards. The legislature proposes listing only drug trafficking and corruption as predicate offenses. What is the problem with this approach?
		\end{block}
		\vspace{0.05cm}
		\stepbox{\scriptsize
			\keypoint{Correct Answer:} A predicate offense is a criminal offense that generates proceeds that are then laundered through the financial system. Under \textbf{FATF Recommendation 3}, countries should apply the ML offense to \textbf{all serious offenses}, or at least adopt a \textbf{list of all designated categories of offenses}. Limiting predicate offenses to only drug trafficking and corruption would:
			\par \textbf{Violate FATF Standards}: FATF requires coverage of all serious offenses.
			\par \textbf{Create Gaps}: Offenses like fraud, human trafficking, and tax evasion would not be predicate offenses.
			\par \textbf{Weaken Enforcement}: Prosecutors could not pursue ML charges for other crimes.
			\par \textbf{Fail EU Standards}: The EU 6AMLD requires 22 specified predicate offenses.
			The FATF-designated categories include: participation in organized crime, terrorism, drug trafficking, corruption, fraud, counterfeiting, human trafficking, and environmental crimes.
		}
		\vspace{0.03cm}
		\tipbox{\scriptsize \textbf{Key Insight:} Under \textit{FATF standards, all serious offenses should be predicate offenses for ML}. The EU 6AMLD expanded to 22 predicate offenses.}
	\end{frame}
	
	% ================================================================
	% SLIDE 14: Q12 — Drug Trafficking (EXPANDED)
	% ================================================================
	\begin{frame}
		\frametitle{\fontsize{11}{13}\selectfont Q12: Drug Trafficking and ML}
		\begin{block}{\fontsize{8}{10}\selectfont Topic: Predicate Crimes — Drug Trafficking}
			\scriptsize
			A financial intelligence unit notes that drug trafficking organizations are increasingly using trade-based money laundering to move proceeds. Why is drug trafficking a persistent source of money laundering?
		\end{block}
		\vspace{0.05cm}
		\stepbox{\scriptsize
			\keypoint{Correct Answer:} Drug trafficking generates enormous \textbf{cash proceeds} that must be laundered to be used in the legitimate economy. Key characteristics:
			\par \textbf{Cash-Intensive}: Drug sales generate physical currency.
			\par \textbf{High Volume}: Large quantities of cash need placement.
			\par \textbf{Cross-Border}: International drug trade requires cross-border money movement.
			\par \textbf{Trade-Based}: Increasing use of trade-based ML to disguise proceeds.
			\par \textbf{Structured}: Use of structuring to avoid reporting thresholds.
			The \textit{cash-intensive nature} of drug trafficking makes it a major source of ML proceeds.
		}
		\vspace{0.03cm}
		\tipbox{\scriptsize \textbf{Key Insight:} \textit{The UN Office on Drugs and Crime estimates that 50\% of global drug trafficking proceeds are laundered through the financial system.}}
	\end{frame}
	
	% ================================================================
	% SLIDE 15: Q13 — Corruption (EXPANDED)
	% ================================================================
	\begin{frame}
		\frametitle{\fontsize{11}{13}\selectfont Q13: Corruption and ML}
		\begin{block}{\fontsize{8}{10}\selectfont Topic: Predicate Crimes — Corruption}
			\scriptsize
			A foreign official receives secret payments in exchange for awarding government contracts. The official uses shell companies and real estate purchases to hide the funds. What is the relationship between corruption and money laundering?
		\end{block}
		\vspace{0.05cm}
		\stepbox{\scriptsize
			\keypoint{Correct Answer:} Corruption proceeds (bribes, embezzlement) must be laundered through the financial system to appear legitimate. The relationship is \textbf{interconnected}:
			\par \textbf{Corruption generates proceeds}: Bribes and embezzled funds are illicit.
			\par \textbf{ML enables corruption}: Laundering allows corrupt officials to use funds without detection.
			\par \textbf{ML hides corruption}: Obscures the connection between officials and illicit funds.
			\par \textbf{Reciprocal Relationship}: Corruption enables ML (by weakening enforcement) and ML enables corruption (by concealing proceeds).
			\textit{FATF recognizes the link} and emphasizes tackling both.
		}
		\vspace{0.03cm}
		\tipbox{\scriptsize \textbf{Key Insight:} \textit{Corruption and money laundering are twin crimes} — one enables the other. The UN Convention against Corruption recognizes this link.}
	\end{frame}
	
	% ================================================================
	% SLIDE 16: Q14 — Human Trafficking (EXPANDED)
	% ================================================================
	\begin{frame}
		\frametitle{\fontsize{11}{13}\selectfont Q14: Human Trafficking and ML}
		\begin{block}{\fontsize{8}{10}\selectfont Topic: Predicate Crimes — Human Trafficking}
			\scriptsize
			An analyst identifies a pattern of small wire transfers from multiple victims' accounts to a single beneficiary, with funds then moved to offshore accounts. This is associated with a human trafficking network. Why is human trafficking associated with money laundering?
		\end{block}
		\vspace{0.05cm}
		\stepbox{\scriptsize
			\keypoint{Correct Answer:} Human trafficking generates significant proceeds that must be laundered through the financial system:
			\par \textbf{High-Profit}: Trafficking is a high-profit, low-risk crime for traffickers.
			\par \textbf{Multiple Revenue Streams}: Exploitation generates ongoing income.
			\par \textbf{Cross-Border}: International trafficking requires cross-border money movement.
			\par \textbf{Small Transactions}: Victims' payments are often small, avoiding thresholds.
			\par \textbf{Diverse Methods}: Proceeds are laundered through wire transfers, crypto, and trade.
			\par \textbf{Integration}: Funds are integrated through real estate, businesses, and investments.
			\textit{FATF recognizes human trafficking as a predicate offense for ML.}
		}
		\vspace{0.03cm}
		\tipbox{\scriptsize \textbf{Key Insight:} \textit{FATF has issued specific guidance on human trafficking and ML}, recognizing the link between exploitation and money laundering.}
	\end{frame}
	
	% ================================================================
	% SLIDE 17: Q15 — ML in Banking (EXPANDED)
	% ================================================================
	\begin{frame}
		\frametitle{\fontsize{11}{13}\selectfont Q15: ML in Banking Sector}
		\begin{block}{\fontsize{8}{10}\selectfont Topic: Manifestation in Banking}
			\scriptsize
			A money launderer deposits \$50,000 in cash into a bank account, immediately wire transfers the funds to three different offshore banks, and then uses the funds to purchase a real estate property. What stages of ML are represented?
		\end{block}
		\vspace{0.05cm}
		\stepbox{\scriptsize
			\keypoint{Correct Answer:} This represents all three stages:
			\par \textbf{Placement}: \$50,000 cash deposit into the bank account.
			\par \textbf{Layering}: Wire transfers to three offshore banks to obscure origin.
			\par \textbf{Integration}: Real estate purchase with the "cleaned" funds.
			Additional banking ML methods include:
			\par \textbf{Structuring}: Breaking deposits into amounts under reporting thresholds.
			\par \textbf{Multiple Accounts}: Using numerous accounts to avoid detection.
			\par \textbf{Shell Companies}: Using corporate entities to hide ownership.
			\par \textbf{Correspondent Banking}: Using foreign correspondent banks to obscure origin.
			\textit{Banking is the most common sector for ML due to ease of moving funds.}
		}
		\vspace{0.03cm}
		\tipbox{\scriptsize \textbf{Key Insight:} \textit{All three stages may occur in the banking sector} — placement through deposits, layering through transfers, integration through investments.}
	\end{frame}
	
	% ================================================================
	% SLIDE 18: Q16 — TF in Banking (EXPANDED)
	% ================================================================
	\begin{frame}
		\frametitle{\fontsize{11}{13}\selectfont Q16: TF in Banking Sector}
		\begin{block}{\fontsize{8}{10}\selectfont Topic: TF in Banking}
			\scriptsize
			A compliance team notices a pattern of small wire transfers (\$200-\$500) from multiple individuals to a single beneficiary in a conflict zone. The transfers reference "charitable donations." What should the team suspect?
		\end{block}
		\vspace{0.05cm}
		\stepbox{\scriptsize
			\keypoint{Correct Answer:} Terrorist financing through small value wire transfers:
			\par \textbf{Small Amounts}: TF often uses small amounts to avoid reporting thresholds (e.g., under \$10,000).
			\par \textbf{Multiple Sources}: Many individuals sending to a single beneficiary.
			\par \textbf{Conflict Zone}: Destination in a conflict zone is a red flag.
			\par \textbf{Charitable Reference}: Use of charitable language to appear legitimate.
			\par \textbf{Multiple Accounts}: Use of multiple accounts to avoid detection.
			The team should:
			\par \textbf{Investigate}: Determine if there is a legitimate charitable purpose.
			\par \textbf{SAR Filing}: File a suspicious activity report if no legitimate basis.
			\par \textbf{CFT Controls}: Ensure CFT controls are in place.
		}
		\vspace{0.03cm}
		\tipbox{\scriptsize \textbf{Key Insight:} \textit{TF often involves small amounts to avoid triggering reporting thresholds.} The \textit{purpose} is key, not the amount.}
	\end{frame}
	
	% ================================================================
	% SLIDE 19: Q17 — Sanctions Evasion (EXPANDED)
	% ================================================================
	\begin{frame}
		\frametitle{\fontsize{11}{13}\selectfont Q17: Sanctions Evasion in Banking}
		\begin{block}{\fontsize{8}{10}\selectfont Topic: Sanctions Evasion}
			\scriptsize
			A bank processes a payment from a company in Country A to a company in Country B, both of which are not sanctioned. However, the beneficial owner of the Country B company is a designated individual. What is this an example of?
		\end{block}
		\vspace{0.05cm}
		\stepbox{\scriptsize
			\keypoint{Correct Answer:} Sanctions evasion using shell companies and third-party intermediaries:
			\par \textbf{Shell Company}: The Country B company is a shell hiding the beneficial owner.
			\par \textbf{Third-Party Intermediary}: The payment is routed through an intermediary to hide the sanctioned party.
			\par \textbf{Complex Structure}: The structure is designed to evade sanctions screening.
			\par \textbf{Ownership/Control Rule}: Under OFAC's 50\% rule, entities owned/controlled by designated persons are also blocked.
			The bank should:
			\par \textbf{Identify UBO}: Conduct beneficial ownership checks.
			\par \textbf{Block Transaction}: Freeze the transaction.
			\par \textbf{Report}: File a sanctions violation report.
			\textit{Sanctions evasion is a strict liability offense} — no intent required.
		}
		\vspace{0.03cm}
		\tipbox{\scriptsize \textbf{Key Insight:} \textit{Sanctions are strict liability} — intent is not required for a violation. The \textit{ownership/control} of entities matters, not just direct transactions.}
	\end{frame}
	
	% ================================================================
	% SLIDE 20: Q18 — Bribery in Banking (EXPANDED)
	% ================================================================
	\begin{frame}
		\frametitle{\fontsize{11}{13}\selectfont Q18: Bribery in Banking}
		\begin{block}{\fontsize{8}{10}\selectfont Topic: Bribery in Banking}
			\scriptsize
			A bank employee accepts payments from a customer in exchange for bypassing KYC controls and approving high-risk accounts. What is this an example of?
		\end{block}
		\vspace{0.05cm}
		\stepbox{\scriptsize
			\keypoint{Correct Answer:} Bribery of bank officials to facilitate financial crime:
			\par \textbf{Bribe}: The employee is receiving something of value.
			\par \textbf{Influence}: The bribe influences the employee's actions.
			\par \textbf{Control Bypass}: KYC controls are being bypassed.
			\par \textbf{Facilitation}: The bribe facilitates other financial crimes (ML, sanctions evasion).
			\par \textbf{Internal Control Failure}: The bank's internal controls have failed.
			The bank should:
			\par \textbf{Investigate}: Determine the extent of the bribery.
			\par \textbf{Terminate}: Fire the employee.
			\par \textbf{Report}: File an internal report and regulatory notification.
			\par \textbf{Remediate}: Strengthen controls and separation of duties.
			\textit{Bribery undermines internal controls and enables other financial crimes.}
		}
		\vspace{0.03cm}
		\tipbox{\scriptsize \textbf{Key Insight:} \textit{Bribery is a gateway crime} — it enables money laundering, sanctions evasion, and other financial crimes.}
	\end{frame}
	
	% ================================================================
	% SLIDE 21: Q19 — Tax Evasion in Banking (EXPANDED)
	% ================================================================
	\begin{frame}
		\frametitle{\fontsize{11}{13}\selectfont Q19: Tax Evasion in Banking}
		\begin{block}{\fontsize{8}{10}\selectfont Topic: Tax Evasion}
			\scriptsize
			A wealthy client maintains multiple offshore accounts, uses shell companies to hold assets, and declares minimal income to tax authorities. The bank's compliance team suspects tax evasion. What are the red flags?
		\end{block}
		\vspace{0.05cm}
		\stepbox{\scriptsize
			\keypoint{Correct Answer:} Tax evasion manifests through:
			\par \textbf{Offshore Accounts}: Accounts in low-tax jurisdictions.
			\par \textbf{Shell Companies}: Entities that conceal ownership.
			\par \textbf{Complex Structures}: Trusts and corporate structures hiding ownership.
			\par \textbf{Inconsistent Declarations}: Assets/income inconsistent with tax declarations.
			\par \textbf{Currency Exchange}: Large currency exchanges or transfers.
			\par \textbf{Frequent Transfers}: Rapid movement of funds through jurisdictions.
			\textit{Tax evasion is a predicate offense for ML in many jurisdictions.}
			The bank should:
			\par \textbf{Conduct EDD}: Enhanced due diligence on the client.
			\par \textbf{Source of Funds}: Verify source of funds and wealth.
			\par \textbf{File SAR}: File a suspicious activity report if evasion is suspected.
			\par \textbf{Comply with CRS/FATCA}: Report to tax authorities as required.
		}
		\vspace{0.03cm}
		\tipbox{\scriptsize \textbf{Key Insight:} \textit{Under the Common Reporting Standard (CRS)}, banks must report offshore accounts to tax authorities. \textit{FATCA} requires US account reporting.}
	\end{frame}
	
	% ================================================================
	% SLIDE 22: Q20 — Fraud in Banking (EXPANDED)
	% ================================================================
	\begin{frame}
		\frametitle{\fontsize{11}{13}\selectfont Q20: Fraud in Banking}
		\begin{block}{\fontsize{8}{10}\selectfont Topic: Fraud in Banking}
			\scriptsize
			A bank identifies a pattern of loan applications using fabricated employment and income documents. The loans are quickly withdrawn and funds transferred overseas. What is this an example of?
		\end{block}
		\vspace{0.05cm}
		\stepbox{\scriptsize
			\keypoint{Correct Answer:} Application fraud and proceeds laundering:
			\par \textbf{Fraud}: Fabricated documents to obtain loans.
			\par \textbf{Placement}: Fraud proceeds deposited into accounts.
			\par \textbf{Integration}: Funds withdrawn and sent overseas.
			\par \textbf{Identity Misuse}: Identity may be stolen or synthetic.
			\par \textbf{Dual Crime}: Fraud generates proceeds that need laundering.
			The bank should:
			\par \textbf{Verify Documents}: Conduct independent verification of income/employment.
			\par \textbf{Identify Patterns}: Look for patterns of similar applications.
			\par \textbf{Block Accounts}: Freeze accounts involved.
			\par \textbf{File Reports}: File fraud reports and SARs.
			\textit{Fraud and money laundering often overlap} — fraud generates proceeds that need laundering.
		}
		\vspace{0.03cm}
		\tipbox{\scriptsize \textbf{Key Insight:} \textit{Fraud and money laundering are often dual crimes} — fraud generates proceeds that must be laundered to be usable.}
	\end{frame}
	
	% ================================================================
	% SLIDE 23: Q21 — Trade Finance Risks (EXPANDED)
	% ================================================================
	\begin{frame}
		\frametitle{\fontsize{11}{13}\selectfont Q21: Trade Finance Risks}
		\begin{block}{\fontsize{8}{10}\selectfont Topic: Banking Product Risks — Trade Finance}
			\scriptsize
			A trade finance analyst reviews an invoice for \$5 million for a shipment of 10,000 tons of steel. The shipment is valued at \$500 per ton. However, the market rate for steel is \$250 per ton. What ML risk does this present?
		\end{block}
		\vspace{0.05cm}
		\stepbox{\scriptsize
			\keypoint{Correct Answer:} Over-invoicing to move value across borders:
			\par \textbf{Over-invoicing}: The invoice is significantly above market rate.
			\par \textbf{Value Transfer}: The extra \$2.5 million (\$250 per ton × 10,000 tons) is being moved as illicit funds.
			\par \textbf{Trade-Based ML}: Using trade to disguise value transfer.
			\par \textbf{Disguise}: The transaction appears as legitimate trade.
			\textit{This is trade-based money laundering (TBML).}
			Additional TBML methods:
			\par \textbf{Under-invoicing}: Invoice below market rate to move money into a country.
			\par \textbf{Misclassification}: Changing goods descriptions to avoid scrutiny.
			\par \textbf{Multiple Invoicing}: Using multiple invoices for the same goods.
			\par \textbf{Over/Under-shipment}: Shipping more/less than declared.
		}
		\vspace{0.03cm}
		\tipbox{\scriptsize \textbf{Key Insight:} \textit{Trade-based ML is difficult to detect} because it appears as legitimate trade. FATF has issued specific guidance on TBML.}
	\end{frame}
	
	% ================================================================
	% SLIDE 24: Q22 — Wealth Management Risks (EXPANDED)
	% ================================================================
	\begin{frame}
		\frametitle{\fontsize{11}{13}\selectfont Q22: Wealth Management Risks}
		\begin{block}{\fontsize{8}{10}\selectfont Topic: Wealth Management}
			\scriptsize
			A wealth management client has \$50 million in assets but the source of wealth is described only as "business success" with no supporting documentation. The client uses multiple trusts and offshore structures. What ML risk does this present?
		\end{block}
		\vspace{0.05cm}
		\stepbox{\scriptsize
			\keypoint{Correct Answer:} Use of offshore structures, trusts, and complex investment vehicles to conceal ownership and source of funds:
			\par \textbf{Complexity}: Trusts and offshore structures hide ownership.
			\par \textbf{Opacity}: Lack of transparency on source of wealth.
			\par \textbf{Vague Descriptions}: "Business success" without documentation.
			\par \textbf{Integration}: Wealth management allows integration of illicit funds.
			\par \textbf{Confidentiality}: Confidentiality can be exploited.
			The institution should:
			\par \textbf{Conduct EDD}: Enhanced due diligence on source of wealth.
			\par \textbf{Verify UBO}: Identify beneficial owners.
			\par \textbf{Monitor Activity}: Monitor for unusual transactions.
			\par \textbf{File SAR}: File suspicious activity report if concerns persist.
			\textit{Wealth management is attractive to launderers due to complexity and confidentiality.}
		}
		\vspace{0.03cm}
		\tipbox{\scriptsize \textbf{Key Insight:} \textit{Wealth management is high-risk} due to the complexity of structures, confidentiality, and ability to integrate large sums.}
	\end{frame}
	
	% ================================================================
	% SLIDE 25: Q23 — Correspondent Banking Risks (EXPANDED)
	% ================================================================
	\begin{frame}
		\frametitle{\fontsize{11}{13}\selectfont Q23: Correspondent Banking Risks}
		\begin{block}{\fontsize{8}{10}\selectfont Topic: Correspondent Banking}
			\scriptsize
			A US bank maintains a correspondent banking relationship with a small bank in a high-risk jurisdiction. The small bank has weak AML controls and is used by clients to access the US financial system. What ML risk does this present?
		\end{block}
		\vspace{0.05cm}
		\stepbox{\scriptsize
			\keypoint{Correct Answer:} The respondent bank's weak AML controls create a channel for illicit funds:
			\par \textbf{Gateway}: The correspondent relationship is a gateway to the US financial system.
			\par \textbf{Weak Controls}: The respondent bank has inadequate AML controls.
			\par \textbf{High-Risk Jurisdiction}: The jurisdiction is high-risk for ML.
			\par \textbf{Shell/Offshore}: Respondent may serve shell or offshore entities.
			\par \textbf{Cross-Border}: Funds can move across borders rapidly.
			\textit{Correspondent banking is a high-risk area for ML.}
			The US bank should:
			\par \textbf{Conduct Due Diligence}: Enhanced due diligence on the respondent.
			\par \textbf{Review Controls}: Assess the respondent's AML controls.
			\par \textbf{Terminate Relationship}: Consider termination if controls are inadequate.
			\par \textbf{File Reports}: File SARs on suspicious activity.
		}
		\vspace{0.03cm}
		\tipbox{\scriptsize \textbf{Key Insight:} \textit{FATF Recommendation 13} requires enhanced due diligence for correspondent banking. \textit{De-risking} is a risk when banks terminate relationships.}
	\end{frame}
	
	% ================================================================
	% SLIDE 26: Q24 — Private Banking Risks (EXPANDED)
	% ================================================================
	\begin{frame}
		\frametitle{\fontsize{11}{13}\selectfont Q24: Private Banking Risks}
		\begin{block}{\fontsize{8}{10}\selectfont Topic: Private Banking}
			\scriptsize
			A private banking client requests a numbered account and asks for "complete discretion" in managing the relationship. The client is a PEP from a high-risk jurisdiction. What risks does this present?
		\end{block}
		\vspace{0.05cm}
		\stepbox{\scriptsize
			\keypoint{Correct Answer:} Private banking high-risk features include:
			\par \textbf{Numbered Accounts}: Accounts that hide ownership.
			\par \textbf{Discretion}: Personalized service can be exploited.
			\par \textbf{PEP}: PEPs are high-risk for corruption.
			\par \textbf{High-Risk Jurisdiction}: Jurisdiction with weak AML controls.
			\par \textbf{High Assets}: Large assets require significant laundering.
			\textit{Private banking is high-risk for ML.}
			The institution should:
			\par \textbf{Conduct EDD}: Enhanced due diligence on the PEP client.
			\par \textbf{Senior Approval}: Senior management approval is required.
			\par \textbf{Source of Funds}: Verify source of funds and wealth.
			\par \textbf{Monitor Activity}: Ongoing monitoring of the relationship.
			\par \textbf{File SAR}: File suspicious activity report if concerns persist.
		}
		\vspace{0.03cm}
		\tipbox{\scriptsize \textbf{Key Insight:} \textit{FATF Recommendation 12} requires enhanced due diligence for PEPs. Private banking is a high-risk area for PEP-related ML.}
	\end{frame}
	
	% ================================================================
	% SLIDE 27: Q25 — DNFBP Risks (EXPANDED)
	% ================================================================
	\begin{frame}
		\frametitle{\fontsize{11}{13}\selectfont Q25: Banking High-Risk Sectors — DNFBPs}
		\begin{block}{\fontsize{8}{10}\selectfont Topic: High-Risk Sectors}
			\scriptsize
			A lawyer handles client funds in an escrow account and facilitates the purchase of real estate using a shell company. The lawyer does not conduct CDD on the beneficial owner. What risk does this present?
		\end{block}
		\vspace{0.05cm}
		\stepbox{\scriptsize
			\keypoint{Correct Answer:} DNFBPs can facilitate money laundering without knowing it:
			\par \textbf{Gatekeeper Role}: Lawyers and accountants can facilitate transactions.
			\par \textbf{Escrow Accounts}: Client funds pass through the lawyer's accounts.
			\par \textbf{Shell Company}: The shell company hides ownership.
			\par \textbf{No CDD}: The lawyer fails to identify the beneficial owner.
			\par \textbf{Real Estate}: Real estate purchases integrate illicit funds.
			\textit{DNFBPs are covered under FATF Recommendation 22.}
			The lawyer should:
			\par \textbf{Conduct CDD}: Identify the client and beneficial owner.
			\par \textbf{File SAR}: File suspicious activity report if concerns persist.
			\par \textbf{Decline Service}: Decline service if CDD cannot be completed.
			\par \textbf{Training}: Ensure training on AML requirements.
		}
		\vspace{0.03cm}
		\tipbox{\scriptsize \textbf{Key Insight:} \textit{DNFBPs are covered under FATF Recommendation 22.} They must conduct CDD and file SARs.}
	\end{frame}
	
	% ================================================================
	% SLIDE 28: Q26 — Casino Risks (EXPANDED)
	% ================================================================
	\begin{frame}
		\frametitle{\fontsize{11}{13}\selectfont Q26: Banking High-Risk Sectors — Casinos}
		\begin{block}{\fontsize{8}{10}\selectfont Topic: Casinos}
			\scriptsize
			A casino customer purchases \$100,000 in chips using cash, plays \$5,000 in games, and cashes out the remaining \$95,000. What ML risk does this represent?
		\end{block}
		\vspace{0.05cm}
		\stepbox{\scriptsize
			\keypoint{Correct Answer:} Casino-based placement and layering:
			\par \textbf{Placement}: \$100,000 cash is introduced into the casino system.
			\par \textbf{Layering}: Minimal play creates the appearance of gambling.
			\par \textbf{Integration}: \$95,000 is cashed out as "gaming winnings."
			\par \textbf{Cash-Intensive}: Casinos handle large amounts of cash.
			\par \textbf{Commingling}: Illicit and legitimate funds are commingled.
			\textit{Casinos are covered under FATF Recommendation 22.}
			The casino should:
			\par \textbf{Conduct CDD}: Identify customers at set thresholds.
			\par \textbf{Monitor}: Monitor transactions for unusual activity.
			\par \textbf{File SAR}: File suspicious activity report for unusual play.
			\par \textbf{Record-Keeping}: Maintain records of transactions.
		}
		\vspace{0.03cm}
		\tipbox{\scriptsize \textbf{Key Insight:} \textit{Casinos can be used for placement and layering.} The "minimal play, full cashout" pattern is a classic red flag.}
	\end{frame}
	
	% ================================================================
	% SLIDE 29: Q27 — Real Estate Risks (EXPANDED)
	% ================================================================
	\begin{frame}
		\frametitle{\fontsize{11}{13}\selectfont Q27: Banking High-Risk Sectors — Real Estate}
		\begin{block}{\fontsize{8}{10}\selectfont Topic: Real Estate}
			\scriptsize
			A foreign client purchases a \$5 million property using a shell company, pays cash, and does not appear to occupy the property. What ML risk does this represent?
		\end{block}
		\vspace{0.05cm}
		\stepbox{\scriptsize
			\keypoint{Correct Answer:} Real estate is attractive for integration:
			\par \textbf{Asset Purchase}: Illicit funds are converted into a legitimate asset.
			\par \textbf{Shell Company}: The shell company hides beneficial ownership.
			\par \textbf{Cash Purchase}: Cash purchase without financing suggests illicit funds.
			\par \textbf{No Occupancy}: The property may be used for value storage, not use.
			\par \textbf{Appreciation}: Real estate may appreciate, increasing value.
			\textit{Real estate is a preferred integration vehicle.}
			The institution should:
			\par \textbf{Identify UBO}: Identify the beneficial owner of the shell company.
			\par \textbf{Verify SOF}: Verify source of funds for the purchase.
			\par \textbf{Conduct EDD}: Enhanced due diligence for the transaction.
			\par \textbf{File SAR}: File suspicious activity report if concerns persist.
		}
		\vspace{0.03cm}
		\tipbox{\scriptsize \textbf{Key Insight:} \textit{Real estate is attractive for integration} — converting illicit funds into legitimate assets. Shell companies are a key red flag.}
	\end{frame}
	
	% ================================================================
	% SLIDE 30: Q28 — Cryptocurrency Risks (EXPANDED)
	% ================================================================
	\begin{frame}
		\frametitle{\fontsize{11}{13}\selectfont Q28: Banking High-Risk Sectors — Cryptocurrency}
		\begin{block}{\fontsize{8}{10}\selectfont Topic: Cryptocurrency}
			\scriptsize
			A cryptocurrency exchange processes rapid trades between Bitcoin and Monero, followed by withdrawal to a privacy wallet. What ML risk does this represent?
		\end{block}
		\vspace{0.05cm}
		\stepbox{\scriptsize
			\keypoint{Correct Answer:} Cryptocurrency exchanges are attractive for ML:
			\par \textbf{Pseudonymity}: Transactions are pseudonymous.
			\par \textbf{Rapid Transfer}: Funds can move rapidly across borders.
			\par \textbf{Privacy Coins}: Monero provides additional anonymity.
			\par \textbf{Privacy Wallets}: Wallets that obscure holdings and transactions.
			\par \textbf{Cross-Border}: Cryptocurrency facilitates cross-border value transfer.
			\textit{VASPs are covered under FATF Recommendation 16.}
			The exchange should:
			\par \textbf{Conduct CDD}: Identify customers (KYC).
			\par \textbf{Monitor}: Monitor for suspicious transactions.
			\par \textbf{Travel Rule}: Comply with the Travel Rule (originator/beneficiary info).
			\par \textbf{File SAR}: File suspicious activity report for suspicious patterns.
		}
		\vspace{0.03cm}
		\tipbox{\scriptsize \textbf{Key Insight:} \textit{VASPs are covered under FATF Recommendation 16 (Travel Rule).} Privacy coins and mixers are challenges for tracing.}
	\end{frame}
	
	% ================================================================
	% SLIDE 31: Q29 — PEP Definition (EXPANDED)
	% ================================================================
	\begin{frame}
		\frametitle{\fontsize{11}{13}\selectfont Q29: PEP Definition}
		\begin{block}{\fontsize{8}{10}\selectfont Topic: PEP and High-Risk Customers}
			\scriptsize
			A bank's compliance team is reviewing a new account for a senior government official from a foreign country. The official is not on any sanctions list but holds a prominent public position. What category of customer is this?
		\end{block}
		\vspace{0.05cm}
		\stepbox{\scriptsize
			\keypoint{Correct Answer:} A foreign Politically Exposed Person (PEP):
			\par \textbf{Definition}: An individual entrusted with a prominent public function.
			\par \textbf{Foreign PEP}: From another country (highest risk).
			\par \textbf{Prominent Position}: Senior government official, diplomat, etc.
			\par \textbf{Higher Risk}: Higher risk of corruption or bribery.
			\par \textbf{FATF Coverage}: Covered under FATF Recommendation 12.
			PEP categories include:
			\par \textbf{Foreign PEPs}: From other countries (highest risk).
			\par \textbf{Domestic PEPs}: From the institution's country (lower risk).
			\par \textbf{International Org PEPs}: In international bodies.
			\par \textbf{Family and Close Associates}: Also covered.
		}
		\vspace{0.03cm}
		\tipbox{\scriptsize \textbf{Key Insight:} \textit{FATF Recommendation 12} requires enhanced due diligence for PEPs. \textit{Foreign PEPs} are the highest risk.}
	\end{frame}
	
	% ================================================================
	% SLIDE 32: Q30 — PEP Risk Factors (EXPANDED)
	% ================================================================
	\begin{frame}
		\frametitle{\fontsize{11}{13}\selectfont Q30: PEP Risk Factors}
		\begin{block}{\fontsize{8}{10}\selectfont Topic: PEP Risk Factors}
			\scriptsize
			A PEP uses a complex structure of shell companies and trusts to hold assets, with a third-party intermediary acting as the registered owner. What is the greatest concern from an AML perspective?
		\end{block}
		\vspace{0.05cm}
		\stepbox{\scriptsize
			\keypoint{Correct Answer:} Indirect ownership intermediaries:
			\par \textbf{Opacity}: Shell companies and trusts hide beneficial ownership.
			\par \textbf{Indirect Ownership}: Third-party intermediaries conceal the PEP's involvement.
			\par \textbf{Offshore Accounts}: Offshore structures avoid scrutiny.
			\par \textbf{Inconsistent Income}: Assets inconsistent with known income.
			\par \textbf{High-Risk}: PEPs using indirect ownership intermediaries present the greatest concern.
			\textit{FATF Recommendation 12 requires enhanced due diligence for PEPs.}
			The institution should:
			\par \textbf{Identify UBO}: Identify the beneficial owner.
			\par \textbf{Verify SOF}: Verify source of funds and wealth.
			\par \textbf{Senior Approval}: Obtain senior management approval.
			\par \textbf{Monitor}: Conduct ongoing monitoring.
		}
		\vspace{0.03cm}
		\tipbox{\scriptsize \textbf{Key Insight:} \textit{PEPs using indirect ownership intermediaries present the greatest concern.} Opacity is the key risk factor.}
	\end{frame}
	
	% ================================================================
	% SLIDE 33: Q31 — High-Risk Customer Indicators (EXPANDED)
	% ================================================================
	\begin{frame}
		\frametitle{\fontsize{11}{13}\selectfont Q31: High-Risk Customer Indicators}
		\begin{block}{\fontsize{8}{10}\selectfont Topic: High-Risk Customers}
			\scriptsize
			A new customer has a complex ownership structure involving multiple jurisdictions, no clear business purpose, and transactions to high-risk jurisdictions. What does this indicate?
		\end{block}
		\vspace{0.05cm}
		\stepbox{\scriptsize
			\keypoint{Correct Answer:} Complex structures + opacity + high-risk jurisdiction = high-risk:
			\par \textbf{Complex Ownership}: Multiple jurisdictions, layers of entities.
			\par \textbf{Opacity}: Beneficial ownership is not clear.
			\par \textbf{High-Risk Jurisdictions}: Transactions involving high-risk countries.
			\par \textbf{No Business Purpose}: Unclear business rationale for the structure.
			\par \textbf{Unusual Patterns}: Transactions inconsistent with the customer's profile.
			\textit{Combined risk factors indicate high-risk customers.}
			The institution should:
			\par \textbf{Conduct EDD}: Enhanced due diligence on the customer.
			\par \textbf{Identify UBO}: Identify beneficial owners.
			\par \textbf{Verify SOF}: Verify source of funds and wealth.
			\par \textbf{Monitor}: Conduct ongoing monitoring.
			\par \textbf{Decline Service}: Consider declining the relationship.
		}
		\vspace{0.03cm}
		\tipbox{\scriptsize \textbf{Key Insight:} \textit{Complexity + Opacity + High-Risk Jurisdiction} = high-risk customer requiring enhanced due diligence.}
	\end{frame}
	
	% ================================================================
	% SLIDE 34: Q32 — PEP EDD Requirements (EXPANDED)
	% ================================================================
	\begin{frame}
		\frametitle{\fontsize{11}{13}\selectfont Q32: PEP EDD Requirements}
		\begin{block}{\fontsize{8}{10}\selectfont Topic: PEP EDD}
			\scriptsize
			A bank is onboarding a foreign PEP with assets of \$20 million. What enhanced due diligence is required under FATF Recommendation 12?
		\end{block}
		\vspace{0.05cm}
		\stepbox{\scriptsize
			\keypoint{Correct Answer:} Enhanced due diligence is mandatory:
			\par \textbf{Source of Funds/Wealth}: Verify source of funds and wealth.
			\par \textbf{Senior Management Approval}: Approval from senior management.
			\par \textbf{Ongoing Monitoring}: Conduct ongoing monitoring of the relationship.
			\par \textbf{Risk Assessment}: Assess the level of risk.
			\par \textbf{Family/Associates}: Identify family members and close associates.
			\textit{EDD for PEPs is mandatory under FATF Recommendation 12.}
			Additional requirements:
			\par \textbf{Politically Exposed Persons}: EDD is required for all PEPs.
			\par \textbf{Enhanced Controls}: Additional controls for high-risk PEPs.
			\par \textbf{Reporting}: File SARs if concerns persist.
		}
		\vspace{0.03cm}
		\tipbox{\scriptsize \textbf{Key Insight:} \textit{FATF Recommendation 12 requires mandatory EDD for PEPs.} It is not optional.}
	\end{frame}
	
	% ================================================================
	% SLIDE 35: Q33 — Securities Sector Risks (EXPANDED)
	% ================================================================
	\begin{frame}
		\frametitle{\fontsize{11}{13}\selectfont Q33: Securities Sector Risks}
		\begin{block}{\fontsize{8}{10}\selectfont Topic: Other Financial Sector Risks}
			\scriptsize
			A securities firm identifies a pattern of suspicious trading with no economic purpose, rapidly crossing between accounts. What ML risk does this represent?
		\end{block}
		\vspace{0.05cm}
		\stepbox{\scriptsize
			\keypoint{Correct Answer:} Market manipulation generates illicit profits:
			\par \textbf{Market Manipulation}: Wash trading, spoofing, insider trading.
			\par \textbf{Illicit Profits}: Manipulation generates illicit proceeds.
			\par \textbf{Integration}: Proceeds can be integrated through investments.
			\par \textbf{Cross-Border}: Securities trading facilitates cross-border value transfer.
			\par \textbf{Complexity}: Securities transactions can be complex and opaque.
			\textit{The securities sector can be used for both generating and laundering illicit funds.}
			The firm should:
			\par \textbf{Monitor Trading}: Implement post-trade surveillance.
			\par \textbf{Identify Patterns}: Look for suspicious patterns.
			\par \textbf{File Reports}: File SARs and market abuse reports.
			\par \textbf{Cooperate}: Cooperate with regulatory investigations.
		}
		\vspace{0.03cm}
		\tipbox{\scriptsize \textbf{Key Insight:} \textit{The securities sector can be used for both generating and laundering illicit funds.} IOSCO has guidance on ML in securities markets.}
	\end{frame}
	
	% ================================================================
	% SLIDE 36: Q34 — Investment Fund Risks (EXPANDED)
	% ================================================================
	\begin{frame}
		\frametitle{\fontsize{11}{13}\selectfont Q34: Investment Fund Risks}
		\begin{block}{\fontsize{8}{10}\selectfont Topic: Investment Funds}
			\scriptsize
			An investment fund accepts a large investment from a shell company in a high-risk jurisdiction, with no clear source of funds. The fund invests in real estate and legitimate businesses. What ML risk does this represent?
		\end{block}
		\vspace{0.05cm}
		\stepbox{\scriptsize
			\keypoint{Correct Answer:} Investment funds can be used for integration:
			\par \textbf{Integration}: Illicit funds are invested in legitimate investments.
			\par \textbf{Shell Company}: The shell company hides beneficial ownership.
			\par \textbf{High-Risk Jurisdiction}: The jurisdiction is high-risk for ML.
			\par \textbf{No SOF}: The source of funds is not verified.
			\par \textbf{Commingling}: Illicit and legitimate investments are commingled.
			\textit{Investment funds can be used for integration of laundered funds.}
			The fund should:
			\par \textbf{Conduct CDD}: Identify investors and beneficial owners.
			\par \textbf{Verify SOF}: Verify source of funds.
			\par \textbf{Monitor}: Monitor for suspicious activity.
			\par \textbf{File SAR}: File suspicious activity report if concerns persist.
		}
		\vspace{0.03cm}
		\tipbox{\scriptsize \textbf{Key Insight:} \textit{Investment funds can be used for integration of laundered funds.} The commingling of illicit and legitimate funds is a key risk.}
	\end{frame}
	
	% ================================================================
	% SLIDE 37: Q35 — Leasing and Financing Risks (EXPANDED)
	% ================================================================
	\begin{frame}
		\frametitle{\fontsize{11}{13}\selectfont Q35: Leasing and Financing Risks}
		\begin{block}{\fontsize{8}{10}\selectfont Topic: Leasing and Financing}
			\scriptsize
			A client leases a luxury vehicle for \$10,000 per month, but the lease is for \$15,000 above market rate. The extra \$5,000 is returned to the client through a third party. What ML risk does this represent?
		\end{block}
		\vspace{0.05cm}
		\stepbox{\scriptsize
			\keypoint{Correct Answer:} Over-valuation in asset financing:
			\par \textbf{Over-Valuation}: Asset value is inflated above market rate.
			\par \textbf{Value Transfer}: The difference is moved to the client.
			\par \textbf{Third-Party Payment}: Payment to a third party obscures the transfer.
			\par \textbf{Integration}: The lease creates the appearance of legitimate value.
			\par \textbf{Complexity}: Leasing arrangements can be complex.
			\textit{Asset financing can be used for over-invoicing and value transfer.}
			The institution should:
			\par \textbf{Verify Value}: Verify the market value of the asset.
			\par \textbf{Conduct CDD}: Identify the client and beneficial owner.
			\par \textbf{File SAR}: File suspicious activity report if concerns persist.
			\par \textbf{Monitor}: Monitor for unusual patterns.
		}
		\vspace{0.03cm}
		\tipbox{\scriptsize \textbf{Key Insight:} \textit{Asset financing can be used for over-invoicing and value transfer.} Over-valuation is a key red flag.}
	\end{frame}
	
	% ================================================================
	% SLIDE 38: Q36 — Pension Fund Risks (EXPANDED)
	% ================================================================
	\begin{frame}
		\frametitle{\fontsize{11}{13}\selectfont Q36: Pension Fund Risks}
		\begin{block}{\fontsize{8}{10}\selectfont Topic: Pension Funds}
			\scriptsize
			A client invests illicit funds in a pension fund, waits several years, and then withdraws the funds as "retirement benefits." What ML risk does this represent?
		\end{block}
		\vspace{0.05cm}
		\stepbox{\scriptsize
			\keypoint{Correct Answer:} Pension funds provide a long-term integration vehicle:
			\par \textbf{Long-Term Integration}: Funds remain in the fund for years.
			\par \textbf{Legitimization}: Withdrawal appears as legitimate retirement benefits.
			\par \textbf{Investment Growth}: Funds may appreciate, increasing value.
			\par \textbf{Difficult to Trace}: The passage of time makes tracing difficult.
			\par \textbf{Low Scrutiny}: Pension funds may receive less scrutiny than other investments.
			\textit{Pension funds provide a long-term integration vehicle for laundered funds.}
			The institution should:
			\par \textbf{Conduct CDD}: Identify investors and beneficial owners.
			\par \textbf{Verify SOF}: Verify source of funds.
			\par \textbf{Monitor}: Monitor for unusual contributions/withdrawals.
			\par \textbf{File SAR}: File suspicious activity report if concerns persist.
		}
		\vspace{0.03cm}
		\tipbox{\scriptsize \textbf{Key Insight:} \textit{Pension funds provide a long-term integration vehicle.} The passage of time makes tracing difficult.}
	\end{frame}
	
	% ================================================================
	% SLIDE 39: Q37 — MSB Risks (EXPANDED)
	% ================================================================
	\begin{frame}
		\frametitle{\fontsize{11}{13}\selectfont Q37: MSB Risks}
		\begin{block}{\fontsize{8}{10}\selectfont Topic: MSB, PSP, E-Commerce}
			\scriptsize
			A Money Services Business (MSB) has agents in multiple high-risk jurisdictions, with no oversight of their AML controls. What risk does this present?
		\end{block}
		\vspace{0.05cm}
		\stepbox{\scriptsize
			\keypoint{Correct Answer:} MSB agent networks influence operational and geographic risk:
			\par \textbf{Agent Network}: Agents in multiple jurisdictions extend the MSB's reach.
			\par \textbf{Operational Risk}: Agents may have weak AML controls.
			\par \textbf{Geographic Risk}: Agents in high-risk jurisdictions increase risk.
			\par \textbf{Cross-Border}: MSBs facilitate cross-border value transfer.
			\par \textbf{Cash-Intensive}: MSBs handle large amounts of cash.
			\textit{MSB agent networks influence both operational and geographic risk exposure.}
			The institution should:
			\par \textbf{Conduct Due Diligence}: Due diligence on agents and jurisdictions.
			\par \textbf{Monitor Agents}: Monitor agent activity for suspicious patterns.
			\par \textbf{Terminate Relationships}: Terminate relationships with weak agents.
			\par \textbf{File SAR}: File suspicious activity report for suspicious activity.
		}
		\vspace{0.03cm}
		\tipbox{\scriptsize \textbf{Key Insight:} \textit{MSB agent networks influence both operational and geographic risk exposure.} Weak agents in high-risk jurisdictions are a key risk.}
	\end{frame}
	
	% ================================================================
	% SLIDE 40: Q38 — PSP Risks (EXPANDED)
	% ================================================================
	\begin{frame}
		\frametitle{\fontsize{11}{13}\selectfont Q38: PSP Risks}
		\begin{block}{\fontsize{8}{10}\selectfont Topic: PSP Risks}
			\scriptsize
			A Payment Service Provider (PSP) processes payments for merchants in high-risk jurisdictions, with no verification of the merchants' legitimacy. What risk does this present?
		\end{block}
		\vspace{0.05cm}
		\stepbox{\scriptsize
			\keypoint{Correct Answer:} PSPs are gateways for fraud and ML:
			\par \textbf{Fraud Facilitation}: PSPs can facilitate fraudulent transactions.
			\par \textbf{Jurisdictional Risk}: Operating in jurisdictions with limited oversight.
			\par \textbf{Merchant Verification}: Merchants may not be legitimate.
			\par \textbf{Cross-Border}: PSPs facilitate cross-border value transfer.
			\par \textbf{High Volume}: PSPs process high volumes of transactions.
			\textit{PSPs are gateways for e-commerce fraud and money laundering.}
			The PSP should:
			\par \textbf{Conduct CDD}: Conduct due diligence on merchants.
			\par \textbf{Monitor}: Monitor for suspicious transactions.
			\par \textbf{File SAR}: File suspicious activity report for suspicious activity.
			\par \textbf{Terminate}: Terminate relationships with high-risk merchants.
		}
		\vspace{0.03cm}
		\tipbox{\scriptsize \textbf{Key Insight:} \textit{PSPs are gateways for e-commerce fraud and money laundering.} Merchant verification is a key control.}
	\end{frame}
	
	% ================================================================
	% SLIDE 41: Q39 — E-Commerce Platform Risks (EXPANDED)
	% ================================================================
	\begin{frame}
		\frametitle{\fontsize{11}{13}\selectfont Q39: E-Commerce Platform Risks}
		\begin{block}{\fontsize{8}{10}\selectfont Topic: E-Commerce}
			\scriptsize
			An e-commerce platform has multiple seller accounts with low sales volume but high transaction frequency, using fake transactions to move value. What ML risk does this represent?
		\end{block}
		\vspace{0.05cm}
		\stepbox{\scriptsize
			\keypoint{Correct Answer:} Trade-based money laundering through e-commerce:
			\par \textbf{Fake Transactions}: Transactions that don't represent real goods/services.
			\par \textbf{Over/Under-Invoicing}: Over/under-invoicing to move value.
			\par \textbf{Fake Products}: Products that don't exist or are misrepresented.
			\par \textbf{Fake Seller Accounts}: Multiple accounts to avoid detection.
			\par \textbf{Cross-Border}: E-commerce facilitates cross-border value transfer.
			\textit{E-commerce platforms can be used for trade-based money laundering.}
			The platform should:
			\par \textbf{Conduct CDD}: Conduct due diligence on sellers.
			\par \textbf{Monitor}: Monitor for suspicious transactions.
			\par \textbf{Sanctions Screening}: Screen sellers against sanctions lists.
			\par \textbf{Tiered Onboarding}: Implement tiered onboarding based on risk.
			\par \textbf{File SAR}: File suspicious activity report for suspicious activity.
		}
		\vspace{0.03cm}
		\tipbox{\scriptsize \textbf{Key Insight:} \textit{E-commerce platforms can be used for trade-based money laundering.} Fake transactions are a key red flag.}
	\end{frame}
	
	% ================================================================
	% SLIDE 42: Q40 — Mitigating E-Commerce Risks (EXPANDED)
	% ================================================================
	\begin{frame}
		\frametitle{\fontsize{11}{13}\selectfont Q40: Mitigating E-Commerce Risks}
		\begin{block}{\fontsize{8}{10}\selectfont Topic: E-Commerce Mitigation}
			\scriptsize
			An e-commerce platform wants to reduce its ML risk while maintaining business growth. What strategies should it implement?
		\end{block}
		\vspace{0.05cm}
		\stepbox{\scriptsize
			\keypoint{Correct Answer:} Tiered onboarding, sanctions screening, and monitoring:
			\par \textbf{Tiered Onboarding}: Different due diligence levels based on projected sales.
			\par \textbf{Sanctions Screening}: Screen all sellers during onboarding and ongoing.
			\par \textbf{Transaction Monitoring}: Monitor for suspicious patterns.
			\par \textbf{Risk Assessment}: Assess and update risk assessments.
			\par \textbf{Training}: Train staff on ML risks and red flags.
			\textit{Tiered onboarding based on risk is an effective control.}
			Additional measures:
			\par \textbf{Adverse Media}: Conduct adverse media checks on sellers.
			\par \textbf{Transaction Limits}: Implement transaction limits based on risk.
			\par \textbf{Reporting}: File SARs for suspicious activity.
			\par \textbf{Cooperation}: Cooperate with law enforcement and FIUs.
		}
		\vspace{0.03cm}
		\tipbox{\scriptsize \textbf{Key Insight:} \textit{Tiered onboarding based on risk is an effective control.} Screening, monitoring, and training are also essential.}
	\end{frame}
	
	% ================================================================
	% SLIDE 43: Q41 — Life Insurance Risks (EXPANDED)
	% ================================================================
	\begin{frame}
		\frametitle{\fontsize{11}{13}\selectfont Q41: Life Insurance Risks}
		\begin{block}{\fontsize{8}{10}\selectfont Topic: Insurance Product Risks}
			\scriptsize
			A client purchases a \$5 million life insurance policy with a single premium payment, then surrenders the policy after two years. What ML risk does this represent?
		\end{block}
		\vspace{0.05cm}
		\stepbox{\scriptsize
			\keypoint{Correct Answer:} Life insurance can be used for placement and integration:
			\par \textbf{Placement}: Large single premium payment introduces illicit funds.
			\par \textbf{Integration}: Surrender returns "cleaned" funds.
			\par \textbf{Wealth Concealment}: Policy conceals the existence of wealth.
			\par \textbf{High-Value}: Single premium policies are high-value.
			\par \textbf{Early Surrender}: Early surrender is a red flag.
			\textit{Life insurance can be used for placement and integration.}
			The insurer should:
			\par \textbf{Conduct CDD}: Identify the policyholder and beneficiary.
			\par \textbf{Verify SOF}: Verify source of funds for the premium.
			\par \textbf{Monitor}: Monitor for early surrender or unusual activity.
			\par \textbf{File SAR}: File suspicious activity report for suspicious patterns.
		}
		\vspace{0.03cm}
		\tipbox{\scriptsize \textbf{Key Insight:} \textit{Life insurance can be used for placement and integration.} Early surrender is a key red flag.}
	\end{frame}
	
	% ================================================================
	% SLIDE 44: Q42 — Property/Casualty Insurance Risks (EXPANDED)
	% ================================================================
	\begin{frame}
		\frametitle{\fontsize{11}{13}\selectfont Q42: Property/Casualty Insurance Risks}
		\begin{block}{\fontsize{8}{10}\selectfont Topic: Property/Casualty}
			\scriptsize
			A client purchases high-value art insurance for a collection, then files a claim for the art's destruction, receiving a large payout. What ML risk does this represent?
		\end{block}
		\vspace{0.05cm}
		\stepbox{\scriptsize
			\keypoint{Correct Answer:} Insurance fraud converts illicit funds into legitimate claim payments:
			\par \textbf{Fraudulent Claim}: False claim for a loss that did not occur.
			\par \textbf{Integration}: Claim payment converts illicit funds into legitimate funds.
			\par \textbf{Over-Valuation}: Art may be over-valued to increase the claim.
			\par \textbf{Arson/Theft}: Loss may be staged to collect the claim.
			\par \textbf{Complexity}: Insurance fraud can be complex.
			\textit{Insurance fraud generates proceeds that need laundering.}
			The insurer should:
			\par \textbf{Verify Claim}: Verify the loss actually occurred.
			\par \textbf{Verify Value}: Verify the value of the insured asset.
			\par \textbf{Investigate}: Investigate suspicious claims.
			\par \textbf{File Reports}: File reports to fraud authorities and regulators.
		}
		\vspace{0.03cm}
		\tipbox{\scriptsize \textbf{Key Insight:} \textit{Insurance fraud generates proceeds that need laundering.} Over-valuation and staged losses are key red flags.}
	\end{frame}
	
	% ================================================================
	% SLIDE 45: Q43 — Reinsurance Risks (EXPANDED)
	% ================================================================
	\begin{frame}
		\frametitle{\fontsize{11}{13}\selectfont Q43: Reinsurance Risks}
		\begin{block}{\fontsize{8}{10}\selectfont Topic: Reinsurance}
			\scriptsize
			A reinsurance arrangement involves multiple layers of risk transfer across jurisdictions, with complex documentation. What ML risk does this present?
		\end{block}
		\vspace{0.05cm}
		\stepbox{\scriptsize
			\keypoint{Correct Answer:} Reinsurance complexity can hide cross-border fund movement:
			\par \textbf{Complexity}: Multiple layers of risk transfer create opacity.
			\par \textbf{Cross-Border}: Reinsurance involves cross-border value transfer.
			\par \textbf{Opaque Documentation}: Documentation may be complex and opaque.
			\par \textbf{Jurisdictional}: Jurisdictional differences may limit oversight.
			\par \textbf{Integration}: Reinsurance can be used to integrate illicit funds.
			\textit{Reinsurance complexity can hide cross-border fund movement.}
			The insurer should:
			\par \textbf{Conduct Due Diligence}: Due diligence on reinsurance partners.
			\par \textbf{Understand Arrangements}: Understand the structure and purpose.
			\par \textbf{Monitor}: Monitor for unusual activity.
			\par \textbf{File SAR}: File suspicious activity report if concerns persist.
		}
		\vspace{0.03cm}
		\tipbox{\scriptsize \textbf{Key Insight:} \textit{Reinsurance complexity can hide cross-border fund movement.} Due diligence and understanding are key.}
	\end{frame}
	
	% ================================================================
	% SLIDE 46: Q44 — Insurance CDD Requirements (EXPANDED)
	% ================================================================
	\begin{frame}
		\frametitle{\fontsize{11}{13}\selectfont Q44: Insurance CDD Requirements}
		\begin{block}{\fontsize{8}{10}\selectfont Topic: Insurance CDD}
			\scriptsize
			An insurance company is required to comply with FATF Recommendation 10. What CDD is required for insurance products?
		\end{block}
		\vspace{0.05cm}
		\stepbox{\scriptsize
			\keypoint{Correct Answer:} Identification of the policyholder and beneficiary:
			\par \textbf{Policyholder}: Identify the person purchasing the policy.
			\par \textbf{Beneficiary}: Identify the person receiving the policy proceeds.
			\par \textbf{CDD Conducted}: CDD must be conducted at the time of application.
			\par \textbf{EDD}: Enhanced due diligence for high-risk policies.
			\par \textbf{Beneficial Owner}: Identify the beneficial owner if a legal entity.
			\textit{Insurance companies are covered under FATF Recommendation 10.}
			Additional requirements:
			\par \textbf{Record-Keeping}: Maintain records of CDD.
			\par \textbf{Monitoring}: Monitor for suspicious activity.
			\par \textbf{Reporting}: File SARs for suspicious activity.
		}
		\vspace{0.03cm}
		\tipbox{\scriptsize \textbf{Key Insight:} \textit{Insurance companies are covered entities under FATF.} CDD must include identification of the policyholder and beneficiary.}
	\end{frame}
	
	% ================================================================
	% SLIDE 47: Q45 — VASP Definition (EXPANDED)
	% ================================================================
	\begin{frame}
		\frametitle{\fontsize{11}{13}\selectfont Q45: VASP Definition}
		\begin{block}{\fontsize{8}{10}\selectfont Topic: VASP and Cryptoassets}
			\scriptsize
			A company operates a cryptocurrency exchange, a wallet service, and a custody service. What category of entity is this under FATF standards?
		\end{block}
		\vspace{0.05cm}
		\stepbox{\scriptsize
			\keypoint{Correct Answer:} A Virtual Asset Service Provider (VASP):
			\par \textbf{Definition}: A person or entity conducting virtual asset exchange, transfer, or custody services for third parties.
			\par \textbf{Exchange}: Facilitating the exchange of virtual assets.
			\par \textbf{Transfer}: Facilitating the transfer of virtual assets.
			\par \textbf{Custody}: Holding virtual assets in custody.
			\par \textbf{FATF Coverage}: VASPs are covered under FATF Recommendation 16.
			\textit{FATF extends AML requirements to VASPs.}
			Activities covered:
			\par \textbf{Exchange-to-Fiat}: Exchange between virtual and fiat currency.
			\par \textbf{Crypto-to-Crypto}: Exchange between virtual assets.
			\par \textbf{Transfer Services}: Transferring virtual assets on behalf of clients.
			\par \textbf{Custody Services}: Safekeeping of virtual assets.
		}
		\vspace{0.03cm}
		\tipbox{\scriptsize \textbf{Key Insight:} \textit{FATF extends AML requirements to VASPs under Recommendation 16.} VASPs must comply with KYC, record-keeping, and reporting requirements.}
	\end{frame}
	
	% ================================================================
	% SLIDE 48: Q46 — Cryptoasset ML Risks (EXPANDED)
	% ================================================================
	\begin{frame}
		\frametitle{\fontsize{11}{13}\selectfont Q46: Cryptoasset ML Risks}
		\begin{block}{\fontsize{8}{10}\selectfont Topic: Cryptoasset Risks}
			\scriptsize
			A cryptocurrency exchange identifies a pattern of rapid transfers between Bitcoin and Monero, followed by withdrawal to a privacy wallet. What ML risk does this represent?
		\end{block}
		\vspace{0.05cm}
		\stepbox{\scriptsize
			\keypoint{Correct Answer:} Pseudonymity, rapid transfer, and privacy coins:
			\par \textbf{Pseudonymity}: Transactions are pseudonymous, not anonymous.
			\par \textbf{Rapid Transfer}: Funds can move rapidly across borders.
			\par \textbf{Privacy Coins}: Monero provides additional anonymity.
			\par \textbf{Mixers}: Mixers (Tornado Cash) obscure transaction trails.
			\par \textbf{Cross-Border}: Cryptocurrency facilitates cross-border value transfer.
			\textit{Privacy coins and mixers are challenges for tracing.}
			The exchange should:
			\par \textbf{Conduct CDD}: Identify customers (KYC).
			\par \textbf{Monitor}: Monitor for suspicious transactions.
			\par \textbf{Travel Rule}: Comply with the Travel Rule.
			\par \textbf{File SAR}: File suspicious activity report for suspicious patterns.
		}
		\vspace{0.03cm}
		\tipbox{\scriptsize \textbf{Key Insight:} \textit{Privacy coins (Monero, Zcash) and mixers (Tornado Cash) are challenges for tracing.} FATF guidance addresses these risks.}
	\end{frame}
	
	% ================================================================
	% SLIDE 49: Q47 — Travel Rule for VASPs (EXPANDED)
	% ================================================================
	\begin{frame}
		\frametitle{\fontsize{11}{13}\selectfont Q47: Travel Rule for VASPs}
		\begin{block}{\fontsize{8}{10}\selectfont Topic: Travel Rule}
			\scriptsize
			A cryptocurrency exchange is implementing the Travel Rule. What information must it collect and share for virtual asset transfers?
		\end{block}
		\vspace{0.05cm}
		\stepbox{\scriptsize
			\keypoint{Correct Answer:} Originator and beneficiary information:
			\par \textbf{Originator Information}: Name, address, and account information.
			\par \textbf{Beneficiary Information}: Name and account information.
			\par \textbf{Transaction Information}: Amount, date, and purpose.
			\par \textbf{Threshold}: Applies to transactions above a threshold.
			\par \textbf{Counterparties}: Must be shared with counterparty VASPs.
			\textit{The Travel Rule applies to VASPs under FATF Recommendation 16.}
			Implementation considerations:
			\par \textbf{Technical}: Technical systems to collect and share information.
			\par \textbf{Privacy}: Compliance with data privacy requirements.
			\par \textbf{Cross-Border}: Applicable to cross-border virtual asset transfers.
		}
		\vspace{0.03cm}
		\tipbox{\scriptsize \textbf{Key Insight:} \textit{The Travel Rule applies to VASPs under FATF Recommendation 16.} It requires collecting and sharing originator and beneficiary information.}
	\end{frame}
	
	% ================================================================
	% SLIDE 50: Q48 — DeFi and ML Risks (EXPANDED)
	% ================================================================
	\begin{frame}
		\frametitle{\fontsize{11}{13}\selectfont Q48: DeFi and ML Risks}
		\begin{block}{\fontsize{8}{10}\selectfont Topic: DeFi Risks}
			\scriptsize
			A Decentralized Finance (DeFi) platform allows users to exchange virtual assets without centralized identity verification. What ML risk does this present?
		\end{block}
		\vspace{0.05cm}
		\stepbox{\scriptsize
			\keypoint{Correct Answer:} Lack of centralized oversight makes tracing difficult:
			\par \textbf{No Central Oversight}: DeFi platforms lack a central authority.
			\par \textbf{No KYC}: Users are not identified or verified.
			\par \textbf{Rapid Transfer}: Virtual assets can be transferred rapidly.
			\par \textbf{Anonymous}: Users are pseudonymous or anonymous.
			\par \textbf{Cross-Border}: DeFi facilitates cross-border value transfer.
			\textit{DeFi platforms can be used for rapid, anonymous value transfer.}
			Challenges:
			\par \textbf{Regulatory Gaps}: DeFi may not be regulated.
			\par \textbf{Enforcement}: Enforcement is difficult.
			\par \textbf{Privacy}: Privacy coins and mixers are common.
			\par \textbf{Tracing}: Tracing is difficult due to anonymity.
		}
		\vspace{0.03cm}
		\tipbox{\scriptsize \textbf{Key Insight:} \textit{DeFi platforms can be used for rapid, anonymous value transfer.} Lack of centralized oversight is a key vulnerability.}
	\end{frame}
	
	% ================================================================
	% SLIDE 51: Q49 — Casino ML Risks (EXPANDED)
	% ================================================================
	\begin{frame}
		\frametitle{\fontsize{11}{13}\selectfont Q49: Casino ML Risks}
		\begin{block}{\fontsize{8}{10}\selectfont Topic: Gaming and Gambling}
			\scriptsize
			A casino customer exchanges \$100,000 in cash for chips, plays a few hands of blackjack, and cashes out the remaining \$97,000. What ML risk does this represent?
		\end{block}
		\vspace{0.05cm}
		\stepbox{\scriptsize
			\keypoint{Correct Answer:} Casino-based placement and layering:
			\par \textbf{Placement}: \$100,000 cash is introduced into the casino system.
			\par \textbf{Layering}: Minimal play creates the appearance of gambling.
			\par \textbf{Integration}: \$97,000 is cashed out as "gaming winnings."
			\par \textbf{Minimal Play}: Minimal play suggests the purpose is not gambling.
			\par \textbf{Cash-Intensive}: Casinos handle large amounts of cash.
			\textit{Casinos can be used for both placement and layering.}
			The casino should:
			\par \textbf{Conduct CDD}: Identify customers at set thresholds.
			\par \textbf{Monitor}: Monitor for unusual play patterns.
			\par \textbf{File SAR}: File suspicious activity report for suspicious patterns.
			\par \textbf{Record-Keeping}: Maintain records of transactions.
		}
		\vspace{0.03cm}
		\tipbox{\scriptsize \textbf{Key Insight:} \textit{Casinos can be used for placement and layering.} Minimal play with full cashout is a classic red flag.}
	\end{frame}
	
	% ================================================================
	% SLIDE 52: Q50 — Online Gambling Risks (EXPANDED)
	% ================================================================
	\begin{frame}
		\frametitle{\fontsize{11}{13}\selectfont Q50: Online Gambling Risks}
		\begin{block}{\fontsize{8}{10}\selectfont Topic: Online Gambling}
			\scriptsize
			An online gambling platform accepts deposits via multiple payment methods and allows withdrawals to different accounts. What ML risk does this present?
		\end{block}
		\vspace{0.05cm}
		\stepbox{\scriptsize
			\keypoint{Correct Answer:} Online gambling is high-risk due to payment complexity:
			\par \textbf{Multiple Payment Methods}: Deposits via credit cards, bank transfers, crypto.
			\par \textbf{Multiple Withdrawal Methods}: Withdrawals to different accounts.
			\par \textbf{Cross-Border}: Online gambling facilitates cross-border value transfer.
			\par \textbf{Anonymous}: Users may be pseudonymous.
			\par \textbf{High Volume}: Online gambling processes high volumes of transactions.
			\textit{Online gambling is high-risk due to cross-border payment complexity.}
			The platform should:
			\par \textbf{Conduct CDD}: Identify users and verify identity.
			\par \textbf{Monitor}: Monitor for suspicious patterns.
			\par \textbf{File SAR}: File suspicious activity report for suspicious activity.
			\par \textbf{Implement Controls}: Implement AML/CFT controls.
		}
		\vspace{0.03cm}
		\tipbox{\scriptsize \textbf{Key Insight:} \textit{Online gambling is high-risk due to cross-border payment complexity.} Multiple payment methods create opacity.}
	\end{frame}
	
	% ================================================================
	% SLIDE 53: Q51 — Sports Betting Risks (EXPANDED)
	% ================================================================
	\begin{frame}
		\frametitle{\fontsize{11}{13}\selectfont Q51: Sports Betting Risks}
		\begin{block}{\fontsize{8}{10}\selectfont Topic: Sports Betting}
			\scriptsize
			A sports betting platform identifies a pattern of large bets on a low-profile match, with the outcomes known in advance. What ML risk does this represent?
		\end{block}
		\vspace{0.05cm}
		\stepbox{\scriptsize
			\keypoint{Correct Answer:} Match-fixing generates proceeds that need laundering:
			\par \textbf{Match-Fixing}: Outcomes are fixed in advance.
			\par \textbf{Insider Betting}: Bets are placed based on inside knowledge.
			\par \textbf{Large Bets}: Large bets on low-profile matches are red flags.
			\par \textbf{Proceeds}: Match-fixing generates illicit proceeds.
			\par \textbf{Integration}: Proceeds must be laundered.
			\textit{Sports betting can be used to generate and launder funds.}
			The platform should:
			\par \textbf{Conduct CDD}: Identify bettors and verify identity.
			\par \textbf{Monitor}: Monitor for suspicious betting patterns.
			\par \textbf{File SAR}: File suspicious activity report for suspicious activity.
			\par \textbf{Cooperate}: Cooperate with sports integrity bodies.
		}
		\vspace{0.03cm}
		\tipbox{\scriptsize \textbf{Key Insight:} \textit{Sports betting can be used to generate and launder funds.} Match-fixing and insider betting are key red flags.}
	\end{frame}
	
	% ================================================================
	% SLIDE 54: Q52 — Gaming CDD Requirements (EXPANDED)
	% ================================================================
	\begin{frame}
		\frametitle{\fontsize{11}{13}\selectfont Q52: Gaming CDD Requirements}
		\begin{block}{\fontsize{8}{10}\selectfont Topic: Gaming CDD}
			\scriptsize
			A casino is required to comply with FATF Recommendations. What CDD is required for gaming establishments?
		\end{block}
		\vspace{0.05cm}
		\stepbox{\scriptsize
			\keypoint{Correct Answer:} Identification at set thresholds, transaction monitoring, and SAR filing:
			\par \textbf{CDD Thresholds}: Identify customers at set thresholds (e.g., \$2,000-\$3,000).
			\par \textbf{Transaction Monitoring}: Monitor for suspicious transactions.
			\par \textbf{SAR Filing}: File SARs for suspicious activity.
			\par \textbf{Record-Keeping}: Maintain records of transactions.
			\par \textbf{EDD}: Enhanced due diligence for high-risk customers.
			\textit{Gaming establishments are covered under FATF.}
			Additional requirements:
			\par \textbf{Cash Reporting}: Report cash transactions above thresholds.
			\par \textbf{AML Program}: Maintain an AML program.
			\par \textbf{Training}: Train employees on AML requirements.
		}
		\vspace{0.03cm}
		\tipbox{\scriptsize \textbf{Key Insight:} \textit{Gaming establishments are covered entities under FATF.} CDD at set thresholds and SAR filing are required.}
	\end{frame}
	
	% ================================================================
	% SLIDE 55: Q53 — Real Estate ML Risks (EXPANDED)
	% ================================================================
	\begin{frame}
		\frametitle{\fontsize{11}{13}\selectfont Q53: Real Estate ML Risks}
		\begin{block}{\fontsize{8}{10}\selectfont Topic: Real Estate}
			\scriptsize
			A foreign buyer purchases a \$10 million property through a shell company, pays cash, and does not appear to occupy the property. What ML risk does this represent?
		\end{block}
		\vspace{0.05cm}
		\stepbox{\scriptsize
			\keypoint{Correct Answer:} Real estate is a preferred integration vehicle:
			\par \textbf{Integration}: Illicit funds are converted into a legitimate asset.
			\par \textbf{Shell Company}: Shell company hides beneficial ownership.
			\par \textbf{Cash Payment}: Cash payment suggests illicit funds.
			\par \textbf{No Occupancy}: Property may be used for value storage.
			\par \textbf{Appreciation}: Real estate may appreciate, increasing value.
			\textit{Real estate is a preferred integration vehicle for laundered funds.}
			The institution should:
			\par \textbf{Identify UBO}: Identify the beneficial owner of the shell company.
			\par \textbf{Verify SOF}: Verify source of funds for the purchase.
			\par \textbf{Conduct EDD}: Enhanced due diligence for the transaction.
			\par \textbf{File SAR}: File suspicious activity report if concerns persist.
		}
		\vspace{0.03cm}
		\tipbox{\scriptsize \textbf{Key Insight:} \textit{Real estate is a preferred integration vehicle.} Shell companies and cash payments are key red flags.}
	\end{frame}
	
	% ================================================================
	% SLIDE 56: Q54 — Real Estate Shell Companies (EXPANDED)
	% ================================================================
	\begin{frame}
		\frametitle{\fontsize{11}{13}\selectfont Q54: Real Estate Shell Companies}
		\begin{block}{\fontsize{8}{10}\selectfont Topic: Shell Companies in Real Estate}
			\scriptsize
			A buyer uses a shell company to purchase a property. The company is registered in a low-tax jurisdiction with no beneficial ownership register. Why are shell companies used in real estate ML?
		\end{block}
		\vspace{0.05cm}
		\stepbox{\scriptsize
			\keypoint{Correct Answer:} Shell companies conceal beneficial ownership:
			\par \textbf{Concealment}: Shell companies hide the true owner.
			\par \textbf{Opacity}: Transparency is lacking.
			\par \textbf{Layering}: Shell companies are a layering tool.
			\par \textbf{Jurisdictional}: Low-tax jurisdictions may have weak transparency.
			\par \textbf{Complexity}: Multiple layers create complexity.
			\textit{Shell companies in real estate are a key red flag.}
			The institution should:
			\par \textbf{Conduct EDD}: Enhanced due diligence on the transaction.
			\par \textbf{Identify UBO}: Identify the beneficial owner.
			\par \textbf{Verify SOF}: Verify source of funds.
			\par \textbf{File SAR}: File suspicious activity report if concerns persist.
		}
		\vspace{0.03cm}
		\tipbox{\scriptsize \textbf{Key Insight:} \textit{Shell companies are a key red flag in real estate ML.} They conceal beneficial ownership.}
	\end{frame}
	
	% ================================================================
	% SLIDE 57: Q55 — Real Estate Over-Valuation (EXPANDED)
	% ================================================================
	\begin{frame}
		\frametitle{\fontsize{11}{13}\selectfont Q55: Real Estate Over-Valuation}
		\begin{block}{\fontsize{8}{10}\selectfont Topic: Over-Valuation}
			\scriptsize
			A property is purchased for \$8 million but is independently valued at \$5 million. The difference of \$3 million is returned to the buyer through a third party. What ML risk does this represent?
		\end{block}
		\vspace{0.05cm}
		\stepbox{\scriptsize
			\keypoint{Correct Answer:} Over-valuation is used for value transfer:
			\par \textbf{Over-Valuation}: The purchase price is above market value.
			\par \textbf{Value Transfer}: The difference is returned to the buyer.
			\par \textbf{Third-Party Payment}: Payment to a third party obscures the transfer.
			\par \textbf{Integration}: The purchase creates a legitimate asset.
			\par \textbf{Appreciation}: The over-valued asset may be sold later.
			\textit{Over-valuation is a form of value transfer.}
			The institution should:
			\par \textbf{Verify Value}: Verify the market value of the property.
			\par \textbf{Investigate}: Investigate the source of the over-valuation.
			\par \textbf{File SAR}: File suspicious activity report if concerns persist.
			\par \textbf{Monitor}: Monitor for similar transactions.
		}
		\vspace{0.03cm}
		\tipbox{\scriptsize \textbf{Key Insight:} \textit{Over-valuation is a form of value transfer.} The difference is returned to the buyer.}
	\end{frame}
	
	% ================================================================
	% SLIDE 58: Q56 — Real Estate EDD Requirements (EXPANDED)
	% ================================================================
	\begin{frame}
		\frametitle{\fontsize{11}{13}\selectfont Q56: Real Estate EDD Requirements}
		\begin{block}{\fontsize{8}{10}\selectfont Topic: Real Estate EDD}
			\scriptsize
			A real estate agent is handling a \$15 million transaction for a foreign buyer using a shell company. What EDD is required?
		\end{block}
		\vspace{0.05cm}
		\stepbox{\scriptsize
			\keypoint{Correct Answer:} Source of funds verification, beneficial ownership identification, and transaction monitoring:
			\par \textbf{SOF Verification}: Verify the source of funds for the purchase.
			\par \textbf{UBO Identification}: Identify the beneficial owner of the shell company.
			\par \textbf{Transaction Monitoring}: Monitor the transaction and subsequent activity.
			\par \textbf{Record-Keeping}: Maintain records of EDD.
			\par \textbf{EDD}: Enhanced due diligence for high-risk transactions.
			\textit{Real estate agents are covered under FATF Recommendation 22.}
			Additional requirements:
			\par \textbf{SAR Filing}: File SARs for suspicious activity.
			\par \textbf{Training}: Train staff on AML requirements.
			\par \textbf{Cooperation}: Cooperate with law enforcement and regulators.
		}
		\vspace{0.03cm}
		\tipbox{\scriptsize \textbf{Key Insight:} \textit{Real estate agents are covered under FATF Recommendation 22.} EDD is required for high-risk transactions.}
	\end{frame}
	
	% ================================================================
	% SLIDE 59: Q57 — Gatekeeper Definition (EXPANDED)
	% ================================================================
	\begin{frame}
		\frametitle{\fontsize{11}{13}\selectfont Q57: Gatekeeper Definition}
		\begin{block}{\fontsize{8}{10}\selectfont Topic: Gatekeeper Risks}
			\scriptsize
			What professional roles are considered gatekeepers in the AML context, and why are they subject to AML requirements?
		\end{block}
		\vspace{0.05cm}
		\stepbox{\scriptsize
			\keypoint{Correct Answer:} Professionals who facilitate financial transactions or provide services that can be exploited for ML:
			\par \textbf{Definition}: Gatekeepers facilitate financial transactions and can be exploited.
			\par \textbf{Roles}: Lawyers, accountants, TCSPs, notaries, real estate agents.
			\par \textbf{Access}: Gatekeepers have access to client funds and information.
			\par \textbf{Concealment}: Gatekeepers can create structures to conceal ownership.
			\par \textbf{Obligations}: Gatekeepers are subject to CDD, reporting, and record-keeping.
			\textit{Gatekeepers are covered under FATF Recommendations 22 and 23.}
			Key obligations:
			\par \textbf{CDD}: Conduct customer due diligence.
			\par \textbf{SAR Filing}: File suspicious activity reports.
			\par \textbf{Record-Keeping}: Maintain records.
			\par \textbf{Training}: Train on AML requirements.
		}
		\vspace{0.03cm}
		\tipbox{\scriptsize \textbf{Key Insight:} \textit{Gatekeepers are covered under FATF Recommendations 22 and 23.} They must conduct CDD and file SARs.}
	\end{frame}
	
	% ================================================================
	% SLIDE 60: Q58 — Lawyer Gatekeeper Risks (EXPANDED)
	% ================================================================
	\begin{frame}
		\frametitle{\fontsize{11}{13}\selectfont Q58: Lawyer Gatekeeper Risks}
		\begin{block}{\fontsize{8}{10}\selectfont Topic: Lawyer Risks}
			\scriptsize
			A lawyer handles client funds in an escrow account, creates a company structure with nominee directors, and facilitates the purchase of real estate. What ML risk does this represent?
		\end{block}
		\vspace{0.05cm}
		\stepbox{\scriptsize
			\keypoint{Correct Answer:} Lawyers can facilitate ML through escrow accounts and corporate structures:
			\par \textbf{Escrow Accounts}: Client funds pass through the lawyer's accounts.
			\par \textbf{Corporate Structures}: Lawyers create companies to conceal ownership.
			\par \textbf{Real Estate}: Lawyers facilitate real estate purchases.
			\par \textbf{Concealment}: Structures are created to conceal ownership.
			\par \textbf{Opacity}: Lack of transparency is a key risk.
			\textit{Lawyers are covered under FATF Recommendation 22.}
			The lawyer should:
			\par \textbf{Conduct CDD}: Identify the client and beneficial owner.
			\par \textbf{Verify SOF}: Verify source of funds.
			\par \textbf{File SAR}: File suspicious activity report if concerns persist.
			\par \textbf{Decline Service}: Decline service if CDD cannot be completed.
		}
		\vspace{0.03cm}
		\tipbox{\scriptsize \textbf{Key Insight:} \textit{Lawyers are covered under FATF Recommendation 22.} Escrow accounts and corporate structures are key risk areas.}
	\end{frame}
	
	% ================================================================
	% SLIDE 61: Q59 — Notary Gatekeeper Risks (EXPANDED)
	% ================================================================
	\begin{frame}
		\frametitle{\fontsize{11}{13}\selectfont Q59: Notary Gatekeeper Risks}
		\begin{block}{\fontsize{8}{10}\selectfont Topic: Notary Risks}
			\scriptsize
			A notary authenticates documents for a real estate transaction involving a shell company, without verifying the beneficial owner. What ML risk does this present?
		\end{block}
		\vspace{0.05cm}
		\stepbox{\scriptsize
			\keypoint{Correct Answer:} Notaries can authenticate documents that conceal ownership:
			\par \textbf{Document Authentication}: Notaries authenticate legal documents.
			\par \textbf{Shell Company}: Documents relate to a shell company.
			\par \textbf{Concealment}: Documents may conceal beneficial ownership.
			\par \textbf{Opacity}: Lack of transparency is a key risk.
			\par \textbf{Real Estate}: Notaries authenticate real estate documents.
			\textit{Notaries are covered under FATF Recommendation 22.}
			The notary should:
			\par \textbf{Conduct CDD}: Identify the parties and beneficial owner.
			\par \textbf{Verify Documents}: Verify the legitimacy of documents.
			\par \textbf{File SAR}: File suspicious activity report if concerns persist.
			\par \textbf{Decline Service}: Decline service if CDD cannot be completed.
		}
		\vspace{0.03cm}
		\tipbox{\scriptsize \textbf{Key Insight:} \textit{Notaries are covered under FATF Recommendation 22.} Document authentication can be exploited to conceal ownership.}
	\end{frame}
	
	% ================================================================
	% SLIDE 62: Q60 — Mitigating Gatekeeper Risks (EXPANDED)
	% ================================================================
	\begin{frame}
		\frametitle{\fontsize{11}{13}\selectfont Q60: Mitigating Gatekeeper Risks}
		\begin{block}{\fontsize{8}{10}\selectfont Topic: Mitigating Gatekeeper Risks}
			\scriptsize
			A law firm wants to mitigate its ML risk while maintaining client relationships. What strategies should it implement?
		\end{block}
		\vspace{0.05cm}
		\stepbox{\scriptsize
			\keypoint{Correct Answer:} Training, client acceptance committees, and standardized due diligence:
			\par \textbf{Training}: Train staff on ML risks and red flags.
			\par \textbf{Client Acceptance Committees}: Review high-risk clients.
			\par \textbf{Standardized DD}: Use standardized due diligence questionnaires.
			\par \textbf{Risk Assessment}: Assess and update risk assessments.
			\par \textbf{Reporting}: File SARs for suspicious activity.
			\textit{Training and client acceptance committees are effective controls.}
			Additional measures:
			\par \textbf{EDD}: Enhanced due diligence for high-risk clients.
			\par \textbf{Monitoring}: Monitor for suspicious activity.
			\par \textbf{Record-Keeping}: Maintain records.
			\par \textbf{Decline Service}: Decline service if CDD cannot be completed.
		}
		\vspace{0.03cm}
		\tipbox{\scriptsize \textbf{Key Insight:} \textit{Training and client acceptance committees are effective controls.} Standardized due diligence is also important.}
	\end{frame}
	
	% ================================================================
	% SLIDE 63: Q61 — TCSP Definition (EXPANDED)
	% ================================================================
	\begin{frame}
		\frametitle{\fontsize{11}{13}\selectfont Q61: TCSP Definition}
		\begin{block}{\fontsize{8}{10}\selectfont Topic: Trust and Company Service Providers}
			\scriptsize
			What services does a Trust and Company Service Provider (TCSP) provide, and why are they subject to AML requirements?
		\end{block}
		\vspace{0.05cm}
		\stepbox{\scriptsize
			\keypoint{Correct Answer:} TCSPs provide company formation, registered office, nominee directors, and trust services:
			\par \textbf{Company Formation}: Creating companies for clients.
			\par \textbf{Registered Office}: Providing a registered office address.
			\par \textbf{Nominee Directors}: Providing nominee directors.
			\par \textbf{Trust Services}: Creating and administering trusts.
			\par \textbf{Concealment}: TCSPs can be used to conceal beneficial ownership.
			\textit{TCSPs are covered under FATF Recommendation 23.}
			Key obligations:
			\par \textbf{CDD}: Conduct customer due diligence.
			\par \textbf{UBO Identification}: Identify beneficial owners.
			\par \textbf{Record-Keeping}: Maintain records.
			\par \textbf{SAR Filing}: File suspicious activity reports.
		}
		\vspace{0.03cm}
		\tipbox{\scriptsize \textbf{Key Insight:} \textit{TCSPs are covered under FATF Recommendation 23.} They can create vehicles for money laundering.}
	\end{frame}
	
	% ================================================================
	% SLIDE 64: Q62 — TCSP ML Risks (EXPANDED)
	% ================================================================
	\begin{frame}
		\frametitle{\fontsize{11}{13}\selectfont Q62: TCSP ML Risks}
		\begin{block}{\fontsize{8}{10}\selectfont Topic: TCSP Risks}
			\scriptsize
			A TCSP forms a company for a client in a high-risk jurisdiction, with nominee directors and a corporate secretary, without identifying the beneficial owner. What ML risk does this represent?
		\end{block}
		\vspace{0.05cm}
		\stepbox{\scriptsize
			\keypoint{Correct Answer:} TCSPs can form companies that conceal beneficial ownership:
			\par \textbf{Company Formation}: Companies are created to conceal ownership.
			\par \textbf{Nominee Directors}: Nominees hide the true owners.
			\par \textbf{Corporate Secretary}: The secretary provides legitimacy.
			\par \textbf{High-Risk Jurisdiction}: Jurisdiction with weak transparency.
			\par \textbf{Opacity}: Lack of transparency is the key risk.
			\textit{TCSPs can create vehicles for money laundering.}
			The TCSP should:
			\par \textbf{Conduct CDD}: Identify the client and beneficial owner.
			\par \textbf{Verify SOF}: Verify source of funds.
			\par \textbf{File SAR}: File suspicious activity report if concerns persist.
			\par \textbf{Decline Service}: Decline service if CDD cannot be completed.
		}
		\vspace{0.03cm}
		\tipbox{\scriptsize \textbf{Key Insight:} \textit{TCSPs can create vehicles for money laundering.} Nominee directors and high-risk jurisdictions are key red flags.}
	\end{frame}
	
	% ================================================================
	% SLIDE 65: Q63 — TCSP Red Flags (EXPANDED)
	% ================================================================
	\begin{frame}
		\frametitle{\fontsize{11}{13}\selectfont Q63: TCSP Red Flags}
		\begin{block}{\fontsize{8}{10}\selectfont Topic: TCSP Red Flags}
			\scriptsize
			A client requests that shareholders be kept confidential, refuses to disclose the beneficial owner, and asks for nominee directors. What should the TCSP do?
		\end{block}
		\vspace{0.05cm}
		\stepbox{\scriptsize
			\keypoint{Correct Answer:} Decline the service due to insufficient CDD:
			\par \textbf{Confidentiality Request}: Request for confidentiality is a red flag.
			\par \textbf{Refusal to Disclose UBO}: Refusal is a red flag.
			\par \textbf{Nominee Directors}: Nominees conceal ownership.
			\par \textbf{High-Risk Jurisdiction}: Jurisdiction may be high-risk.
			\par \textbf{No Business Purpose}: Unclear business purpose is a red flag.
			\textit{TCSPs should decline clients who refuse UBO disclosure.}
			The TCSP should:
			\par \textbf{Conduct CDD}: Conduct enhanced due diligence.
			\par \textbf{Identify UBO}: Attempt to identify the beneficial owner.
			\par \textbf{File SAR}: File suspicious activity report if concerns persist.
			\par \textbf{Decline Service}: Decline the service if CDD cannot be completed.
		}
		\vspace{0.03cm}
		\tipbox{\scriptsize \textbf{Key Insight:} \textit{TCSPs should decline clients who refuse UBO disclosure.} Confidentiality requests are red flags.}
	\end{frame}
	
	% ================================================================
	% SLIDE 66: Q64 — TCSP CDD Requirements (EXPANDED)
	% ================================================================
	\begin{frame}
		\frametitle{\fontsize{11}{13}\selectfont Q64: TCSP CDD Requirements}
		\begin{block}{\fontsize{8}{10}\selectfont Topic: TCSP CDD}
			\scriptsize
			A TCSP is required to comply with FATF Recommendations. What CDD is required?
		\end{block}
		\vspace{0.05cm}
		\stepbox{\scriptsize
			\keypoint{Correct Answer:} Identify the client and beneficial owner, understand the purpose, and monitor:
			\par \textbf{Client Identification}: Identify the client (KYC).
			\par \textbf{UBO Identification}: Identify the beneficial owner.
			\par \textbf{Purpose Understanding}: Understand the business purpose.
			\par \textbf{Monitoring}: Monitor the relationship.
			\par \textbf{Record-Keeping}: Maintain records.
			\textit{CDD for TCSPs is mandatory under FATF.}
			Additional requirements:
			\par \textbf{EDD}: Enhanced due diligence for high-risk clients.
			\par \textbf{SAR Filing}: File suspicious activity reports.
			\par \textbf{Training}: Train staff on AML requirements.
		}
		\vspace{0.03cm}
		\tipbox{\scriptsize \textbf{Key Insight:} \textit{CDD for TCSPs is mandatory under FATF.} Identification of UBO is essential.}
	\end{frame}
	
	% ================================================================
	% SLIDE 67: Q65 — Accountant Gatekeeper Risks (EXPANDED)
	% ================================================================
	\begin{frame}
		\frametitle{\fontsize{11}{13}\selectfont Q65: Accountant Gatekeeper Risks}
		\begin{block}{\fontsize{8}{10}\selectfont Topic: Accountant Risks}
			\scriptsize
			An accountant prepares accounts for a client with complex structures and no clear source of funds. What ML risk does this represent?
		\end{block}
		\vspace{0.05cm}
		\stepbox{\scriptsize
			\keypoint{Correct Answer:} Accountants can prepare accounts that conceal the source of funds:
			\par \textbf{Concealment}: Accounts can conceal illicit activity.
			\par \textbf{Tax Evasion}: Accounts can facilitate tax evasion.
			\par \textbf{Complexity}: Complex structures create opacity.
			\par \textbf{No SOF}: No clear source of funds is a red flag.
			\par \textbf{Opacity}: Lack of transparency is the key risk.
			\textit{Accountants can be used to conceal illicit activity.}
			The accountant should:
			\par \textbf{Conduct CDD}: Identify the client and beneficial owner.
			\par \textbf{Verify SOF}: Verify source of funds.
			\par \textbf{File SAR}: File suspicious activity report if concerns persist.
			\par \textbf{Decline Service}: Decline service if CDD cannot be completed.
		}
		\vspace{0.03cm}
		\tipbox{\scriptsize \textbf{Key Insight:} \textit{Accountants can be used to conceal illicit activity.} Tax evasion and opacity are key risks.}
	\end{frame}
	
	% ================================================================
	% SLIDE 68: Q66 — Accountant Red Flags (EXPANDED)
	% ================================================================
	\begin{frame}
		\frametitle{\fontsize{11}{13}\selectfont Q66: Accountant Red Flags}
		\begin{block}{\fontsize{8}{10}\selectfont Topic: Accountant Red Flags}
			\scriptsize
			A client refuses to provide source of funds, has complex structures without clear business purpose, and asks for creative accounting. What should the accountant do?
		\end{block}
		\vspace{0.05cm}
		\stepbox{\scriptsize
			\keypoint{Correct Answer:} Decline the service and consider SAR filing:
			\par \textbf{Refusal of SOF}: Client refuses to provide source of funds.
			\par \textbf{No Business Purpose}: Complex structures without clear purpose.
			\par \textbf{Creative Accounting}: Requests for creative accounting are red flags.
			\par \textbf{High-Risk}: High-risk jurisdiction or industry.
			\par \textbf{Opacity}: Lack of transparency is a red flag.
			\textit{Accountants must question unusual client requests.}
			The accountant should:
			\par \textbf{Conduct CDD}: Conduct enhanced due diligence.
			\par \textbf{File SAR}: File suspicious activity report if concerns persist.
			\par \textbf{Decline Service}: Decline the service.
			\par \textbf{Report}: Report to the relevant authorities.
		}
		\vspace{0.03cm}
		\tipbox{\scriptsize \textbf{Key Insight:} \textit{Accountants must question unusual client requests.} Red flags include refusal to provide SOF and requests for creative accounting.}
	\end{frame}
	
	% ================================================================
	% SLIDE 69: Q67 — Accountant ML Facilitation (EXPANDED)
	% ================================================================
	\begin{frame}
		\frametitle{\fontsize{11}{13}\selectfont Q67: Accountant ML Facilitation}
		\begin{block}{\fontsize{8}{10}\selectfont Topic: Accountant Facilitation}
			\scriptsize
			An accountant creates false invoices, overstates expenses, and prepares accounts that conceal beneficial ownership. How does this facilitate ML?
		\end{block}
		\vspace{0.05cm}
		\stepbox{\scriptsize
			\keypoint{Correct Answer:} Accountants can be willing or unwilling facilitators:
			\par \textbf{False Records}: False invoices and records conceal activity.
			\par \textbf{Overstated Expenses}: Overstating expenses to move value.
			\par \textbf{Concealed UBO}: Concealing beneficial ownership.
			\par \textbf{Tax Evasion}: Facilitating tax evasion.
			\par \textbf{Integration}: Helping to integrate illicit funds.
			\textit{Accountants can be willing or unwilling facilitators of ML.}
			The accountant should:
			\par \textbf{Verify Records}: Verify the accuracy of records.
			\par \textbf{Question Red Flags}: Question suspicious requests.
			\par \textbf{File SAR}: File suspicious activity report if concerns persist.
			\par \textbf{Decline Service}: Decline service if concerns persist.
		}
		\vspace{0.03cm}
		\tipbox{\scriptsize \textbf{Key Insight:} \textit{Accountants can be willing or unwilling facilitators of ML.} False records and concealed UBO are key red flags.}
	\end{frame}
	
	% ================================================================
	% SLIDE 70: Q68 — Accountant Responsibilities (EXPANDED)
	% ================================================================
	\begin{frame}
		\frametitle{\fontsize{11}{13}\selectfont Q68: Accountant Responsibilities}
		\begin{block}{\fontsize{8}{10}\selectfont Topic: Accountant Responsibilities}
			\scriptsize
			What are accountants' AML responsibilities under FATF Recommendation 22?
		\end{block}
		\vspace{0.05cm}
		\stepbox{\scriptsize
			\keypoint{Correct Answer:} CDD, record-keeping, and SAR filing:
			\par \textbf{CDD}: Conduct customer due diligence.
			\par \textbf{Record-Keeping}: Maintain records of CDD and transactions.
			\par \textbf{SAR Filing}: File suspicious activity reports.
			\par \textbf{Training}: Train on AML requirements.
			\par \textbf{EDD}: Enhanced due diligence for high-risk clients.
			\textit{Accountants are covered under FATF Recommendation 22.}
			Additional responsibilities:
			\par \textbf{Client Acceptance}: Develop client acceptance policies.
			\par \textbf{Risk Assessment}: Assess and update risk assessments.
			\par \textbf{Cooperation}: Cooperate with law enforcement and regulators.
			\par \textbf{Internal Controls}: Maintain internal controls.
		}
		\vspace{0.03cm}
		\tipbox{\scriptsize \textbf{Key Insight:} \textit{Accountants are covered under FATF Recommendation 22.} They must conduct CDD and file SARs.}
	\end{frame}
	
	% ================================================================
	% SLIDE 71: Q69 — Free-Trade Zone Risks (EXPANDED)
	% ================================================================
	\begin{frame}
		\frametitle{\fontsize{11}{13}\selectfont Q69: Free-Trade Zone Risks}
		\begin{block}{\fontsize{8}{10}\selectfont Topic: Other High-Risk Sectors}
			\scriptsize
			A Free-Trade Zone (FTZ) has lax customs inspection protocols and limited regulatory oversight. What ML risk does this present?
		\end{block}
		\vspace{0.05cm}
		\stepbox{\scriptsize
			\keypoint{Correct Answer:} FTZs enable trade-based ML due to limited inspection:
			\par \textbf{Lack of Inspection}: Limited inspection protocols.
			\par \textbf{Transshipment}: Goods are transshipped through the FTZ.
			\par \textbf{Misclassification}: Goods can be misclassified.
			\par \textbf{Value Transfer}: Over/under-invoicing is common.
			\par \textbf{Opacity}: Lack of transparency is a key risk.
			\textit{FTZs are high-risk due to limited oversight.}
			The institution should:
			\par \textbf{Conduct CDD}: Identify parties and verify documentation.
			\par \textbf{Verify Documents}: Verify shipping and customs documents.
			\par \textbf{Monitor}: Monitor for suspicious patterns.
			\par \textbf{File SAR}: File suspicious activity report if concerns persist.
		}
		\vspace{0.03cm}
		\tipbox{\scriptsize \textbf{Key Insight:} \textit{FTZs are high-risk due to limited oversight.} Lack of inspection is the key vulnerability.}
	\end{frame}
	
	% ================================================================
	% SLIDE 72: Q70 — Precious Metals and Stones Risks (EXPANDED)
	% ================================================================
	\begin{frame}
		\frametitle{\fontsize{11}{13}\selectfont Q70: Precious Metals and Stones Risks}
		\begin{block}{\fontsize{8}{10}\selectfont Topic: Precious Metals/Stones}
			\scriptsize
			A customer purchases \$500,000 in gold bullion with cash, transports it to another country, and sells it there. What ML risk does this represent?
		\end{block}
		\vspace{0.05cm}
		\stepbox{\scriptsize
			\keypoint{Correct Answer:} Precious metals are high-value, portable, and easily concealed:
			\par \textbf{High-Value}: Precious metals can store large value in small volume.
			\par \textbf{Portable}: Easy to transport and conceal.
			\par \textbf{Cross-Border}: Easy to move across borders.
			\par \textbf{Value Transfer}: Used for value transfer across borders.
			\par \textbf{Integration}: Can be sold to integrate illicit funds.
			\textit{Precious metals are covered under FATF Recommendation 23.}
			The institution should:
			\par \textbf{Conduct CDD}: Identify the customer and verify identity.
			\par \textbf{Verify SOF}: Verify source of funds for the purchase.
			\par \textbf{File SAR}: File suspicious activity report if concerns persist.
			\par \textbf{Monitor}: Monitor for suspicious patterns.
		}
		\vspace{0.03cm}
		\tipbox{\scriptsize \textbf{Key Insight:} \textit{Precious metals are covered under FATF Recommendation 23.} They are high-value and portable.}
	\end{frame}
	
	% ================================================================
	% SLIDE 73: Q71 — Art and Antiquities Risks (EXPANDED)
	% ================================================================
	\begin{frame}
		\frametitle{\fontsize{11}{13}\selectfont Q71: Art and Antiquities Risks}
		\begin{block}{\fontsize{8}{10}\selectfont Topic: Art and Antiquities}
			\scriptsize
			A high-value artwork is sold through a private sale with no public record, using subjective valuation to justify a high price. What ML risk does this represent?
		\end{block}
		\vspace{0.05cm}
		\stepbox{\scriptsize
			\keypoint{Correct Answer:} Art is attractive for value transfer and integration:
			\par \textbf{Subjective Valuation}: Art values are subjective and can be manipulated.
			\par \textbf{Over/Under-Invoicing}: Over/under-invoicing to move value.
			\par \textbf{Anonymous Sales}: Private sales can be anonymous.
			\par \textbf{Cross-Border}: Art can be moved across borders.
			\par \textbf{Integration}: Art sales can integrate illicit funds.
			\textit{Art is attractive for value transfer and integration.}
			The institution should:
			\par \textbf{Conduct CDD}: Identify the parties and beneficial owners.
			\par \textbf{Verify Value}: Verify the value of the artwork.
			\par \textbf{Verify SOF}: Verify source of funds.
			\par \textbf{File SAR}: File suspicious activity report if concerns persist.
		}
		\vspace{0.03cm}
		\tipbox{\scriptsize \textbf{Key Insight:} \textit{Art is attractive for value transfer and integration.} Subjective valuation and anonymous sales are key risks.}
	\end{frame}
	
	% ================================================================
	% SLIDE 74: Q72 — Shell Companies and Trusts (EXPANDED)
	% ================================================================
	\begin{frame}
		\frametitle{\fontsize{11}{13}\selectfont Q72: Shell Companies and Trusts}
		\begin{block}{\fontsize{8}{10}\selectfont Topic: Shell Companies/Trusts}
			\scriptsize
			A client uses a shell company and a trust structure to hold assets across multiple jurisdictions, with no clear business purpose. What ML risk does this represent?
		\end{block}
		\vspace{0.05cm}
		\stepbox{\scriptsize
			\keypoint{Correct Answer:} Shell companies and trusts conceal beneficial ownership:
			\par \textbf{Concealment}: Shell companies and trusts hide true ownership.
			\par \textbf{Layering}: They are used for layering illicit funds.
			\par \textbf{Opacity}: Lack of transparency is the key risk.
			\par \textbf{Cross-Border}: Structures span multiple jurisdictions.
			\par \textbf{No Business Purpose}: No clear business purpose is a red flag.
			\textit{Shell companies and trusts are layering tools.}
			The institution should:
			\par \textbf{Conduct CDD}: Identify the client and beneficial owner.
			\par \textbf{Verify SOF}: Verify source of funds and wealth.
			\par \textbf{Understand Purpose}: Understand the business purpose.
			\par \textbf{File SAR}: File suspicious activity report if concerns persist.
		}
		\vspace{0.03cm}
		\tipbox{\scriptsize \textbf{Key Insight:} \textit{Shell companies and trusts are layering tools.} They conceal beneficial ownership and create opacity.}
	\end{frame}
	
	% ================================================================
	% SLIDE 75: SUMMARY TABLE
	% ================================================================
	\begin{frame}
		\frametitle{\fontsize{11}{13}\selectfont SUMMARY: Key Principles — Domain A}
		\scriptsize
		\begin{block}{\fontsize{9}{11}\selectfont Understanding Financial Crime Risks:}
			\vspace{0.05cm}
			\scriptsize
			\begin{longtable}{|p{0.30\textwidth}|p{0.65\textwidth}|}
				\hline
				\keypoint{Concept} & \keypoint{Application / Key Insight} \\
				\hline
				Money Laundering & 3 stages: Placement → Layering → Integration \\
				\hline
				Terrorist Financing & Funds may be legitimate; \textit{intent} is key \\
				\hline
				Sanctions & Strict liability; OFAC, UN, EU regimes \\
				\hline
				Fraud & Identity theft vs. synthetic identity fraud \\
				\hline
				Anti-Bribery & UK Act prohibits facilitation; FCPA allows \\
				\hline
				Tax Evasion & Illegal concealment vs. legal avoidance \\
				\hline
				ML Consequences & Economic, reputational, social impacts \\
				\hline
				Predicate Crimes & Drug trafficking, corruption, human trafficking \\
				\hline
				Banking Products & Trade finance, wealth management, correspondent, private banking \\
				\hline
				High-Risk Sectors & DNFBPs, casinos, real estate, crypto \\
				\hline
				PEPs & Foreign, domestic, int'l org; \textit{EDD mandatory} \\
				\hline
				MSB/PSP/E-Commerce & Agent networks, fraud facilitation, TBML \\
				\hline
				Insurance & Life, property/casualty, reinsurance risks \\
				\hline
				VASP/Crypto & Travel Rule, privacy coins, DeFi risks \\
				\hline
				Gaming/Gambling & Casinos, online gambling, sports betting \\
				\hline
				Real Estate & Shell companies, over-valuation, cash purchases \\
				\hline
				Gatekeepers & Lawyers, accountants, TCSPs, notaries \\
				\hline
				TCSPs & Company formation, nominee directors, trusts \\
				\hline
				FTZs & Limited inspection; trade-based ML \\
				\hline
				Precious Metals/Art & Subjective valuation; value transfer \\
				\hline
			\end{longtable}
			\vspace{0.05cm}
		\end{block}
		\vspace{0.05cm}
		\tipbox{\scriptsize \textbf{FINAL LESSON:} Domain A tests understanding of \textbf{how financial crime works} — the definitions, methods, and sector-specific risks. Understand the \textit{methods} and the \textit{answers} become obvious.}
	\end{frame}
	
	% ================================================================
	% SLIDE 76: EXAM TIPS
	% ================================================================
	\begin{frame}
		\frametitle{\fontsize{11}{13}\selectfont EXAM TIPS — Domain A}
		\scriptsize
		\begin{block}{\fontsize{9}{11}\selectfont Common Traps to Avoid:}
			\vspace{0.05cm}
			\scriptsize
			\begin{itemize}
				\item \highlight{TF vs ML}: TF funds can be \textbf{legitimate}; ML funds are \textbf{always illicit}.
				\item \highlight{Tax Evasion vs Avoidance}: Evasion is \textbf{illegal}; avoidance is \textbf{legal}.
				\item \highlight{Facilitation Payments}: Allowed under FCPA; \textbf{prohibited} under UK Bribery Act.
				\item \highlight{Placement vs Layering}: Placement = \textbf{entering} the system; Layering = \textbf{moving} to obscure origin.
				\item \highlight{Integration}: Converting illicit funds into \textbf{legitimate} assets.
				\item \highlight{PEP EDD}: EDD is \textbf{mandatory} under FATF Recommendation 12.
				\item \highlight{VASP Travel Rule}: VASPs must share originator/beneficiary information.
				\item \highlight{FTZ Risk}: \textbf{Lack of inspection} is the key vulnerability.
				\item \highlight{Gatekeepers}: Lawyers, accountants, TCSPs are \textbf{covered entities}.
				\item \highlight{Sanctions}: Strict liability — \textbf{no intent required}.
			\end{itemize}
			\vspace{0.05cm}
		\end{block}
		\vspace{0.05cm}
		\tipbox{\scriptsize \textbf{Key Insight:} Focus on understanding the \textit{methods} criminals use in each sector. The questions test your ability to \textit{apply} concepts, not just memorize definitions.}
	\end{frame}
	
	% ================================================================
	% SLIDE 77: THANK YOU
	% ================================================================
	\begin{frame}
		\centering
		\vspace{2.5cm}
		{\fontsize{16}{19}\selectfont \textbf{Thank You}} \\[0.2cm]
		{\fontsize{11}{13}\selectfont Keep learning. Keep growing.} \\[0.15cm]
		\rule{3cm}{0.5pt} \\[0.15cm]
		{\fontsize{8}{10}\selectfont \textit{Understanding financial crime risks is the foundation of effective AML/CFT compliance.}}
		\vspace{0.3cm}
		{\fontsize{7}{9}\selectfont \textcolor{darkgray}{End of Domain A Review — 77 Slides}}
	\end{frame}
	
\end{document}